\section{Genealogical closure and cardinality of the generated continuum}
\label{sec:cierre-cardinal-nativo}

This section brings together, in a contiguous proof, extension, double
projection, and arbitrary depth; it then constructs the semantic closure,
the uniform lift, the genealogical machinery, and the cardinal closure. The
set-theoretic hierarchy is identified only at the end of the construction and
cannot be read as its causal origin. The causal dependence is:

\[
 \mathrm{APP}\longrightarrow\mathrm{TRIT}\longrightarrow\mathrm{TPK}
 \longrightarrow X_\infty^{\rm enr}
 \longrightarrow\mathcal C_{\rm cont}^{\rm disc}
 \longrightarrow m_\infty
 \longrightarrow\mathfrak G_{\rm HMT}(h,m_\infty).
\tag{C.32}\label{eq:ch-cadena-causal}
\]

First, the state, its extensions, the continuum, and the genealogical ledger
are constructed. Then \textbf{that produced genealogy} is extended to all its
hereditary objects by an explicit semantic rule. The subsequent identification
with a known set-theoretic hierarchy selects neither the seeds, the operators,
the genealogical ledger, the line, nor its projections. It does, however,
permit the exact logical-cardinal strength of the adopted semantic rule to be
assessed with complete transparency.

\subsection{Extension of states and transfinite semantic extension}

Let

\[
 \Sigma_{\rm HMT}
 =\bigl(\dot R_{\rm APP},\dot R_{\rm TRIT},\dot R_{\rm TPK},
 \dot R_{U},\dot R_{\Gamma_9},\dot R_{\rm Ext},\dot R_\rho,
 \dot R_{\pi},\dot R_{\varphi},\dot R_e,\dot R_{\alpha},
 \dot R_{\rm ar},\dot R_{\rm exc}\bigr)
\tag{C.33}\label{eq:ch-signatura}
\]

be the countable \textbf{relational} signature whose effective code is \(h\).
Each \(\dot R_S\) is the graph of the operator \(S\) on the codes where it is
defined; no external function is assumed to be total within each level. The
graphs \(\dot R_{\pi},\dot R_{\varphi},\dot R_e,\dot R_{\alpha}\) are the
output readers constructed by coinductive generation and dodecaphase closure;
they do not take the conventional values of the constants as arguments. Their
presence in \(\Sigma_{\rm HMT}\) ensures that the closure preserves the full
state: it simultaneously integrates the routes that express value and the
incidence that extends toward the exceptional realization.

At each horizon, the state code records the APP--TRIT--TPK coordinates, the
route, sheet, carry, boundary, and memory; at its representation stages it also
records \(P_N,\Phi_N,E_N,A_N\), where \(A_N\) is the prefix of
\(\alpha_{\rm HMT}\) produced by the same state. These coordinates do not
select the branch from outside: they are part of the history being encoded.
Once a primitive-recursive numbering of the finite coordinates of a state is
fixed, every inverse history
\(x_\infty=(x_N)_{N\geq1}\in X_\infty^{\rm enr}\) determines the code
\[
 \operatorname{code}(x_\infty)
 :=\{\langle N,j,c\rangle:(x_N)_j\text{ has code }c\}\subseteq\omega.
\]
The coherence equations
\(\rho_{N+1,N}(x_{N+1})=x_N\) make it possible to decode the history uniquely
from that real. We write \(m_\infty\) for the code of the realized history and
\(m_\infty(n)\) for its characteristic function under the fixed numbering.
The operator code \(h\) and this complete ledger are combined, by means of a
primitive-recursive pairing function, into the real

\[
 u_{h,m}:=
 \{\langle0,\langle n,h(n)\rangle\rangle:n<\omega\}
 \cup
 \{\langle1,\langle n,m_\infty(n)\rangle\rangle:n<\omega\}.
\tag{C.33a}\label{eq:ch-codigo-conjunto}
\]

The two primitive projections recover \(h\) and \(m_\infty\) from
\(u_{h,m}\).

The first extension is the metric and operator extension of the finite states.
For \(N\ge1\), let \(\widehat{\mathcal X}^{\rm enr,raw}_{6N}\) be the set of
raw enriched states of length \(6N\), and define the surviving subsystem
set-theoretically by
\[
\begin{aligned}
 X_N
 &:=\widehat{\mathcal X}^{\rm enr,surv}_{6N}\\
 &:=\Bigl\{x\in\widehat{\mathcal X}^{\rm enr,raw}_{6N}:\\[-.2ex]
 &\qquad\forall M>N\ \exists z\in
   \widehat{\mathcal X}^{\rm enr,raw}_{6M}\;\text{ such that}\\[-.2ex]
 &\qquad\operatorname{tr}_{6M,6N}(z)=x\ \wedge\
   z\text{ satisfies every intermediate TPK transition}\Bigr\}.
\end{aligned}
\]
and set
\[
 \rho_{N+1,N}:=\operatorname{tr}_{6N+6,6N}\!\upharpoonright X_{N+1},
\]
and
\[
 \operatorname{Ext}_N(x):=
 \{y\in X_{N+1}:(x,y)\in
 \operatorname{Ext}_{6N+6,6N}\}.
\]
The root level is included without any implicit exception:
\[
 X_0:=\{\ast\},\qquad
 \rho_{1,0}:X_1\to X_0\text{ is constant},\qquad
 \operatorname{Ext}_0(\ast):=X_1.
\]
The superscript \(\mathrm{surv}\) denotes the pruned tree of states possessing
compatible TPK descendants at every finite horizon. The TPK extension law
established in the preceding sections proves that the emitted seeds belong to
this tree; finite branching and Kőnig's lemma turn compatibility at every
horizon into an infinite branch. In particular,

\[
 \forall N\ge0\;\forall x\in X_N\quad
 \varnothing\ne\operatorname{Ext}_N(x)
 \subseteq\rho_{N+1,N}^{-1}(x).
\tag{C.34}\label{eq:ch-extension-no-vacia}
\]

By the same finitely branching tree lemma, every surviving seed has a
compatible extension

\[
 x_\infty=(x_N)_{N\ge1}\in
 X_\infty^{\rm enr}:=\varprojlim_N(X_N,\rho_{N+1,N}).
\tag{C.35}\label{eq:ch-limite-inverso}
\]

The second extension acts \textbf{afterward} on the complete code of that
realized history. It is not inferred from the mere existence of the finite
operator \(\operatorname{Ext}_N\): at set-theoretic scale it formalizes the
HMT principle that the only objects in the integral domain are those having a
genealogical derivation from the state and its rules. The transfinite semantic
extension is:

\[
 \begin{aligned}
  \mathfrak G_0(h,m_\infty)
  &:=\operatorname{TC}\bigl(\{\omega,u_{h,m},h,m_\infty\}\bigr),\\
  \mathfrak G_{\beta+1}(h,m_\infty)
  &:=\operatorname{Def}_{\Sigma_{\rm HMT}}
       \bigl(\mathfrak G_\beta(h,m_\infty)\bigr),\\
  \mathfrak G_\lambda(h,m_\infty)
  &:=\bigcup_{\beta<\lambda}\mathfrak G_\beta(h,m_\infty)
       \quad(\lambda\text{ limit}),\\
  \mathfrak G_{\rm HMT}(h,m_\infty)
  &:=\bigcup_{\beta\in\operatorname{Ord}}
       \mathfrak G_\beta(h,m_\infty).
 \end{aligned}
\tag{C.36}\label{eq:ch-clausura-transfinita}
\]

The Continuum Hypothesis and a target bijection are not introduced among the
conditions defining the successor operator; the operator is given
prospectively by

\[
 \operatorname{Def}_{\Sigma_{\rm HMT}}(M)
 :=M\cup
 \Bigl\{
   \{x\in M:\langle M,\in,
       (\dot R_S\cap M^{k_S})_{S\in\Sigma_{\rm HMT}}\rangle
       \models\varphi(x,\bar p)\}:
   \ulcorner\varphi\urcorner\in\omega,
   \ \bar p\in M^{<\omega}
 \Bigr\}.
\tag{C.37}\label{eq:ch-operador-def}
\]

That is, it simultaneously expresses all subsets definable by a finite formula
of the HMT relational language and parameters already produced by the
genealogy. It consults neither a target real value, a conventional constant, a
target cardinal, nor an external model. Formula~\eqref{eq:ch-operador-def}
defines the transfinite closure over the previously constructed seeds,
operators, enriched state, ledger, and double projection.

\begin{lemma}[Joint coding of the HMT program and the realized ledger]
\label{thm:ch-6}
The pair formed by the operator code \(h\) and the ledger \(m_\infty\) is
encoded by a single real \(u_{h,m}\subseteq\omega\), and every formula of
\(\Sigma_{\rm HMT}\) is uniformly translated into a formula of the language
\(\{\in,u_{h,m}\}\).

\end{lemma}
\begin{proof}
Both \(h\) and \(m_\infty\) are sequences of natural numbers: \(h\) encodes
the signature, the finite tables, and the operator programs; \(m_\infty\)
encodes one realized history, not the complete history space.
Formula~\eqref{eq:ch-codigo-conjunto} and its two decoders prove the first
assertion. Each symbol in~\eqref{eq:ch-signatura} denotes its \textbf{graph
relation encoded in \(h\)}, never the external graph of all its possible
evaluations. APP, TRIT, \(U_t\), \(\Gamma_9\), the restrictions \(\rho\), and
the finite extensions have primitive-recursive graphs on integer codes and,
therefore, \(\Delta_0\) formulas absolute between transitive levels. The two
representations are expressed as graph relations through the formulas

\[
 \begin{aligned}
 \dot R_{\rm ar}(x,y)
 &\;\Longleftrightarrow\;
 \forall k\in\omega\;\exists N\;\forall n\ge N\quad
 |\mu_n(x)-y|<2^{-k},\\
 \dot R_{\rm exc}(x,c,z)
 &\;\Longleftrightarrow\;
 \forall N\in\omega\quad z\!\upharpoonright N
   =E_N(x\!\upharpoonright N,c\!\upharpoonright N),
 \end{aligned}
\tag{C.37c$_0$}\label{eq:ch-grafos-publicacion}
\]

where \(\mu_n\) and \(E_N\) are the rational/finite functions whose codes
occur in \(h\). Every atomic formula containing an HMT symbol is inductively
replaced by its graph formula with parameter \(u_{h,m}\). Connectives and
quantifiers are preserved. This induction on syntactic complexity produces a
translation \(\varphi\mapsto\varphi^*\) such that, for every transitive level
\(M\) of~\eqref{eq:ch-clausura-transfinita} containing the parameters,

\[
 \langle M,\in,(\dot R_S\cap M^{k_S})_S\rangle
   \models\varphi(\bar a)
 \quad\Longleftrightarrow\quad
 M\models\varphi^*(\bar a;u_{h,m}).
\tag{C.37c}\label{eq:ch-eliminacion-uniforme}
\]

Thus, \eqref{eq:ch-operador-def} cannot consult an oracle, a target decimal
table, or an uncoded relation that introduces uncountably many objects in a
single step.
\end{proof}

HMT semantics is now defined \textbf{independently of identifying it with
\(\mathfrak G\)}. Let \(\mathcal T_0^{\rm HMT}\) be the set of canonical names
of the elements of
\(\mathfrak G_0(h,m_\infty)=
\operatorname{TC}(\{\omega,u_{h,m},h,m_\infty\})\). In particular, the
initial level contains the realized ledger \(m_\infty\), not every history
in \(X_\infty^{\rm enr}\) as primitive parameters. The admissible terms are
the well-founded trees obtained by

\[
 \begin{aligned}
 \mathcal T^{\rm HMT}_{\beta+1}
 &:={\mathcal T}^{\rm HMT}_\beta\cup
 \{\langle\beta,\ulcorner\varphi\urcorner,
       t_0,\ldots,t_{k-1}\rangle:
       t_i\in\mathcal T^{\rm HMT}_\beta\},\\
 \mathcal T^{\rm HMT}_\lambda&:=
       \bigcup_{\beta<\lambda}\mathcal T^{\rm HMT}_\beta,
 \qquad
 \mathcal T^{\rm HMT}:=
       \bigcup_{\beta\in\operatorname{Ord}}\mathcal T^{\rm HMT}_\beta.
 \end{aligned}
\tag{C.37a}\label{eq:ch-terminos}
\]

The value of a successor term is the subset of the preceding level defined by
\(\varphi\) in the induced relational structure, with the values of
\(t_0,\ldots,t_{k-1}\) as parameters. Before any set-theoretic comparison,
define

\[
 \operatorname{Ob}_{\rm sem}^{\rm HMT}(h,m_\infty)
 :=\{\llbracket t\rrbracket:t\in\mathcal T^{\rm HMT}\}.
\tag{C.37b}\label{eq:ch-objetos-semanticos}
\]

This is the rule of \textbf{generative exhaustivity} adopted by HMT: every
object in the integral domain must possess one of these derivation trees, and
every admissible tree expresses an object. It is not a consequence of finite
branching. The Continuum Hypothesis is not introduced into this rule; it is
obtained through the proof developed in the article. Neither an enumeration of
the reals nor a target bijection is introduced as data.

\begin{theorem}[Equivalence between term semantics and the genealogical hierarchy]
\label{thm:ch-6a}
The evaluation of the preceding terms satisfies

\[
 \operatorname{Ob}_{\rm sem}^{\rm HMT}(h,m_\infty)
 =\mathfrak G_{\rm HMT}(h,m_\infty).
\tag{C.37d}\label{eq:ch-equivalencia-semantica}
\]

\end{theorem}
\begin{proof}
By transfinite induction. At the initial level, the two classes coincide. If
every value of a term in \(\mathcal T^{\rm HMT}_\beta\) belongs to
\(\mathfrak G_\beta\), the semantics of a successor term is one of the subsets
in~\eqref{eq:ch-operador-def}. Conversely, every subset in
~\eqref{eq:ch-operador-def} is given by a formula code and a finite tuple of
parameters; by the induction hypothesis those parameters possess terms, and
~\eqref{eq:ch-terminos} forms the term that denotes it. At a limit, both
constructions take the same union.
\end{proof}

\Needspace{9\baselineskip}
The state encoder is compatible with truncations and sends every TPK extension
to a successor term. Let
\(\mathcal T_N^{\rm st}\subset\mathcal T^{\rm HMT}\) be the subclass of terms
encoding complete states of depth \(N\), and let
\(\tau_{N+1,N}:\mathcal T_{N+1}^{\rm st}\to\mathcal T_N^{\rm st}\) delete the
last block together with the corresponding truncation of the genealogical
ledger, carry, and memory. If \(J_N:X_N\to\mathcal T_N^{\rm st}\) is the
complete encoder, compatibility with \(\operatorname{Ext}_N\) proves the
well-typed identity

\[
 \forall x\in X_N\ \forall y\in\operatorname{Ext}_N(x)\qquad
 \tau_{N+1,N}\bigl(J_{N+1}(y)\bigr)=J_N(x).
\tag{C.37e$_0$}\label{eq:ch-naturalidad-codificador}
\]

Define within the semantics
\[
 \operatorname{Succ}^{\rm TPK}_{N+1}(J_N(x)):=
 \{t\in\mathcal T_{N+1}^{\rm st}:\exists y\in X_{N+1}\,
   (y\in\operatorname{Ext}_N(x)\ \wedge\ t=J_{N+1}(y))\}.
\]
Choose \(\beta\) such that \(\mathfrak G_\beta\) contains \(x\), \(X_N\),
\(X_{N+1}\), the finite graphs of
\(\operatorname{Ext}_N,J_N,J_{N+1}\), and their codes in \(h\). Then the
primitive-recursive formula for the graph of \(\operatorname{Ext}_N\) gives

\[
 J_{N+1}[\operatorname{Ext}_N(x)]
 =\operatorname{Succ}^{\rm TPK}_{N+1}(J_N(x)),
 \qquad
 J_{N+1}[\operatorname{Ext}_N(x)]
 \in\mathfrak G_{\beta+1},
\tag{C.37e}\label{eq:ch-sucesores-semanticos}
\]

and both sides are formed by the graph formula of the same TPK extension.
Equations~\eqref{eq:ch-naturalidad-codificador}--\allowbreak\eqref{eq:ch-sucesores-semanticos}
are the explicit semantic link: the subclass of successor terms is exactly the
image of the admissible extensions, and its truncation commutes. The complete
fiber of \(\rho\) is not identified with \(\operatorname{Ext}_N\), nor is
\(\operatorname{Def}\) identified with TPK merely by nomenclature.

\begin{lemma}[Transitivity, axioms, and well-ordering of the genealogical model]
\label{thm:ch-6b}
\(\mathfrak G_{\rm HMT}(h,m_\infty)\) is a transitive class, contains all
ordinals, satisfies Zermelo--Fraenkel set theory (ZF), and possesses a
definable global well-order; it therefore satisfies Choice and all axioms used
in the subsequent cardinal comparison.

\end{lemma}
\begin{proof}
Abbreviate \(B=\mathfrak G_0\) and \(u=u_{h,m}\). Induction on
~\eqref{eq:ch-clausura-transfinita} proves simultaneously that every level is
transitive and that \(\mathfrak G_\alpha\in\mathfrak G_{\alpha+1}\); if all
ordinals smaller than \(\alpha\) are present, their set is definable at the
corresponding level and expresses \(\alpha\). Hence the union contains all
ordinals. Extensionality and Foundation are inherited from the ambient
universe; Empty Set and Infinity belong to \(B\). Once a level containing the
parameters is chosen, the formulas defining \(\{a,b\}\), \(\bigcup a\), and
subsets by Separation express them at a successor level.

For Collection--Replacement, let \(F\) be a functional relation definable in
\(\mathfrak G\) on the set \(a\). The reflection scheme relative to the
definable hierarchy
\((\mathfrak G_\xi,\in,u_{h,m})_{\xi\in\mathrm{Ord}}\), applied in the ambient
metatheory to the finite list of formulas defining \(F\), produces a level
\(\mathfrak G_\theta\) containing \(a\) and the parameters and correct for
those formulas. The ranks of the witnesses in \(F[a]\) are then bounded by
\(\theta\); Separation over \(\mathfrak G_\theta\) forms the image in
\(\mathfrak G_{\theta+1}\). For Power Set, the internal set to be constructed
is not presupposed. Given \(a\in\mathfrak G_\gamma\), consider in the ambient
metatheory the already existing set \(\mathcal P^V(a)\). For every
\(b\in\mathcal P^V(a)\cap\mathfrak G\), let \(r(b)\) be its first
genealogical stage. Since \(\mathfrak G\) is a definable class, Separation in
the ambient metatheory first forms the set
\(\mathcal P^V(a)\cap\mathfrak G\). Replacement \textbf{in the ambient
metatheory}, applied to that set, bounds the ordinals \(r(b)\) by some
\(\theta\). Consequently,
\[
 \mathcal P^{\mathfrak G}(a)
 =\{b\in\mathfrak G_\theta:b\subseteq a\},
\]
and the right-hand side is expressed by Separation in
\(\mathfrak G_{\theta+1}\). It does not introduce into \(\mathfrak G\) the
subsets external to the semantics~\eqref{eq:ch-objetos-semanticos}. This is the
noncircular rank-bounding argument.

Finally, each object is assigned its first level of birth and its least
\textbf{derivation name}: a well-founded, finitely branching derivation tree
(and, by~\eqref{eq:ch-terminos}, finite for each individual name), labeled by
a level ordinal, a formula code, and the tuple of names of its parameters.
These names are ordered by transfinite recursion: first by the ordinal label,
then by the formula code, and then lexicographically by the already ordered
parameter names. The least name of each value defines a class relation
\(<_{\mathfrak G}\) that compares all objects. It is a definable global
well-order and, by taking the least element of every nonempty set, proves
Choice.
\end{proof}

\subsection{TPK lift of the ternary fiber and complete cylinder partition}
\label{subsec:elevacion-catalogal-ch}

The fiber of \(729=3^6\) refinement subcylinders does not arise from the transfer quotient
\(\mathcal K_{\rm ph}\). That operator is a subsequent representation of reduced
memory and cannot generate histories. The covering is constructed directly
in the system \((H_k,\rho_k,\operatorname{Ext}_k)\) through the three
one-turn extensions of the nonadic holonomy.

\subsubsection{Hensel inversion in the arithmetic tower and realization by histories}
\label{sec:ch-capacidad-hensel}

Before lifting the branches to enriched histories, it is useful to recover the
finite argument determining the capacity of the arithmetic tower. The
complete proof for the matrix's unfiltered transducer is
included; its domain is distinguished from the surviving tree.

\begin{proposition}[Conjugacy of the unfiltered Hensel tower]
\label{prop:ch-hensel-bruto}
Let \(p\) be prime, \(d\geq1\) and \(A\in\operatorname{GL}_d(\mathbb Z)\).
For row vectors, taking residue representatives in
\(\{0,\ldots,p-1\}^d\), define
\[
 x_tA=u_t+px_{t+1},\qquad u_t=x_tA\bmod p.
\]
The map
\[
 \Psi_T:(\mathbb Z/p^T\mathbb Z)^d\longrightarrow
             (\mathbb F_p^d)^T,\qquad
 x_0\longmapsto(u_0,\ldots,u_{T-1})
\]
is bijective, with inverse
\[
 x_0\equiv\sum_{t=0}^{T-1}p^tu_tA^{-(t+1)}\pmod{p^T}.
\]
The bijections commute with reductions and give a homeomorphism
\(\Psi:\mathbb Z_p^d\to(\mathbb F_p^d)^{\mathbb N}\), where
\(\mathbb Z_p:=\varprojlim_T\mathbb Z/p^T\mathbb Z\). The update
\(x_t\mapsto x_{t+1}\) corresponds to deleting the first block.
\end{proposition}

\begin{proof}
The integer inverse of \(A\) exists because its determinant is \(\pm1\).
Iterating \(x_t=(u_t+px_{t+1})A^{-1}\) gives
\[
 x_0=\sum_{t=0}^{T-1}p^tu_tA^{-(t+1)}+p^Tx_TA^{-T}.
\]
Modulo \(p^T\), a word determines a unique initial class. Given an
arbitrary word, setting \(x_T=0\) and reconstructing backwards produces a
class emitting it, proving surjectivity. The same formula
proves compatibility under truncation. Passage to the inverse limit gives the
homeomorphism and the relation to the tape shift.
\end{proof}

For \(p=3\) and \(d=6\), the block alphabet has \(3^6=729\)
elements. This arithmetic capacity is not identified by definition with
the realizability of every block after every surviving prefix. The
link to the \(\operatorname{Ext}\) edges is precisely the role
of Lemma~\ref{lem:elevacion-catalogal-nonadica}. The two
operations distinguished by the matrix are thus maintained: representation in the raw tower and
realization under the APP--TRIT--TPK compatibilities.


\subsubsection{Edge-by-edge realization of the three nonadic extensions}
\label{subsubsec:tres-vueltas-ext-ch}

This subsection specifies the autonomous model generated by the TPK extension
relation entering the cardinal proof. The existence of an edge is not inferred
from a coincidence of types: its initial and final states, APP relation, fiber
transport, and unary update of the ledger are constructed. The construction
uses exclusively the two APP sheets, the TRIT selector, and the TPK
composition \(\mathsf{Sel}\)--\(\mathsf{Tra}\)--\(\operatorname{Upd}\).

\paragraph{Complete cells and three balanced reductions.}
In the chart \(D_9=\{1,\ldots,9\}\), we always retain the complete APP cell
\[
 a_{p,q}:=(p,q;R^\Sigma_{p,q},q^\Sigma_{p,q},
                 R^\Pi_{p,q},q^\Pi_{p,q}),
\]
with
\[
 p+q=R^\Sigma_{p,q}+9q^\Sigma_{p,q},\qquad
 pq=R^\Pi_{p,q}+9q^\Pi_{p,q}.
\]
The active sheet is recorded in a separate coordinate; \(a_{p,q}\) is not
split into a purported additive cell and a separate multiplicative cell. For
\(r\in D_9\), let
\begin{equation}
 a(r):=
 \begin{cases}
  a_{1,9},&r=1,\\
  a_{1,r-1},&2\le r\le9,
 \end{cases}
 \qquad
 \tau(r):=M_{\rm ph}(r),
 \qquad
 m(r):=\tau(r)\in\{-1,0,+1\}.
 \label{eq:celula-completa-fase-ch}
\end{equation}
The additive residue of \(a(r)\) supplies an APP representative of phase
here; this does not make \(\Sigma\) the universally active sheet. The dynamic
selector activates \(\Sigma\) for \(r\in\{1,4,7\}\), activates \(\Pi\) for
\(r\in\{2,5,8\}\), and preserves both sheets without moving the cursor for
\(r\in\{3,6,9\}\). In all three cases,
\(R^\Sigma,q^\Sigma,R^\Pi,q^\Pi\) are transported jointly. We have
\[
 (m(1),\ldots,m(9))=(1,-1,0,1,-1,0,1,-1,0).
\]
The three pairs ordered by the nonadic base are
\begin{equation}
 \mathfrak p_{+1}:=(1,4),\qquad
 \mathfrak p_{-1}:=(2,5),\qquad
 \mathfrak p_{0}:=(3,6).
 \label{eq:tres-pares-internos-ch}
\end{equation}
These pairs are the normalized section taking the first two occurrences of
each TRIT type under the shift \(r\mapsto r+3\); the third phase block
\(7,8,9\) retains the turn-closure control. Both members belong to the same
fiber of \(M_{\rm ph}\). Without introducing a target coefficient, the
balanced reduction from the TRIT chapter gives
\begin{equation}
 \begin{array}{c|c|c|c}
 \varepsilon&m(\mathfrak p_\varepsilon^{(1)})
                  +m(\mathfrak p_\varepsilon^{(2)})
   &r_3\bigl(2\varepsilon\bigr)&
   \chi(\varepsilon,\varepsilon)\\ \hline
 +1&1+1=(-1)+3\cdot1&-1&+1\\
 0&0+0=0+3\cdot0&0&0\\
 -1&(-1)+(-1)=1+3\cdot(-1)&+1&-1
 \end{array}
 \label{eq:tres-carry-derivados-ch}
\end{equation}
Therefore,
\begin{equation}
 \chi(\varepsilon,\varepsilon)
 =\frac{2\varepsilon-r_3(2\varepsilon)}3=\varepsilon.
 \label{eq:carry-no-parametrico-ch}
\end{equation}
The branch sign is thus the output of the balanced cocycle applied to two TRIT
emissions produced by APP. Each pair supplies one of the three carry outputs.
For clarity, we subsequently index the pair by
\(\varepsilon(\mathfrak p):=
\chi(m(\mathfrak p^{(1)}),m(\mathfrak p^{(2)}))\); the symbol
\(\varepsilon\) does not supply a free carry to \(\operatorname{Upd}^{\rm cal}\).

The phase architecture, balanced lift, and memory transport are realized here
through twenty-seven local incidences and their extension graphs. The
normalized section fixes representatives of the fibers of \(M_{\rm ph}\); the
covering to be proved uses the existence of these representatives and remains
invariant when they are replaced by another compatible section.

\paragraph{Typed data of the nine edges.}
Let \(K_{\rm Carr}:=\langle\kappa_{\rm Carr}\rangle_{\rm Ab}\) be the carry
module. The actions are typed by
\[
 H_{\rm Carr}:\mathscr P_{\rm TPK}^{\rm op}
   \longrightarrow\operatorname{End}(K_{\rm Carr}),\qquad
 Y_{\rm Carr}:\mathscr P_{\rm TPK}
   \longrightarrow\operatorname{End}(K_{\rm Carr}).
\]
In this generated model, take the trivial actions
\(H_{\rm Carr}(e)=Y_{\rm Carr}(e)=\operatorname{id}_{K_{\rm Carr}}\): the
active sheet and the route are retained as coordinates of the ledger itself
and do not act on the free carry summand. Also let
\((\mathsf M_d,\jmath_{d',d})\) be the graded direct system of memory. To
avoid confusing two representatives identified by the colimit, the state also
retains the degree as a coordinate, and the labeled carrier
\[
 \widetilde{\mathsf M}:=
 \bigsqcup_{d\ge1}\bigl(\{d\}\times\mathsf M_d\bigr),
 \qquad
 \operatorname{gr}_{\mathsf M}(d,\mu)=d.
\]
The representation in the limit memory is given by
\[
\begin{aligned}
 \mathsf M&:=\varinjlim_d\mathsf M_d,&
 \jmath_d&:\mathsf M_d\longrightarrow\mathsf M,\\
 u&:=\jmath_d(u_d),&
 \iota_{\mathsf M}:\mathbb Z&\longrightarrow\mathsf M,
 \qquad n\longmapsto nu .
\end{aligned}
\]
Here \(\jmath_d\) is the canonical colimit map, and the \(u_d\) form a
compatible collection. Its free component permits the reader \(\deg_u(nu)=n\).
The configuration, phase, memory, carry, route, and survival projections will
be typed on the calibrated carrier constructed recursively below; no ambient
space will be imported for them. For every constructed edge \(e\), the
calibrated component
\(\operatorname{Upd}^{\rm cal}_e:=\mathsf{Mem}^{\rm cal}_e\) will denote the
third map of the factorization; its law is fixed field by field below and is
not presupposed as the restriction of an ambient relation.

Fix the initial depth \(k_0=1\), a complete configuration
\(q_\star\in\mathcal Q_{\rm obs}\) of phase one, and
\[
 q_{n,r}:=F_{\rm obs}^{9n+r-1}(q_\star),\qquad
 q_{n,10}:=q_{n+1,1}.
\]
Because the observable calendar is periodic, its value \(q_{n,r}\) does not
by itself distinguish occurrences separated by twelve turns. The calibrated
model therefore uses the labeled lift
\[
 \widetilde q_{n,r}:=(q_{n,r},n,r),\qquad
 \widetilde q_{n,10}:=\widetilde q_{n+1,1},\qquad
 \pi_{\rm obs}(\widetilde q_{n,r})=q_{n,r}.
\]
The lifted dynamics
\(\widetilde F_{\rm obs}(\widetilde q_{n,r})
=\widetilde q_{n,r+1}\) projects to \(F_{\rm obs}\); the label does not alter the physical phase, but makes its
levels of occurrence distinct. The copy of the cell at each level is labeled by
\[
 a_{n,r}:=(n,r;a(r)),\qquad
 (n,r)^+:=
 \begin{cases}
  (n,r+1),&1\le r<9,\\
  (n+1,1),&r=9.
 \end{cases}
\]
The label only prevents two occurrences of the same cell from being
identified; the two Euclidean divisions and their four coordinates are those
of \(a(r)\). For turn \(n\ge0\) and phase \(r\in D_9\), the symbol
\[
 e_{\varepsilon,n,r}:\widetilde q_{n,r}
 \longrightarrow\widetilde q_{n,r+1},
 \qquad r^+:=1+(r\bmod9),
\]
occupies level \(d_{n,r}:=k_n+r-1\), where \(k_n:=k_0+9n\). Its local ledger
contains the complete cell \(a(r)\), the active sheet determined by
\(M_{\rm ph}(r)\), the type \(\tau(r)\), the membership mark in
\(\mathfrak p_\varepsilon\), and, upon reaching the second member of that
pair, the certified identity \eqref{eq:tres-carry-derivados-ch}. We define its
operators solely from that ledger:
\begin{align}
 \mathsf C_{\mathsf M}(e_{\varepsilon,n,r})
   &=\jmath_{d_{n,r}+1,d_{n,r}},&
 \kappa_{\mathsf M}(e_{\varepsilon,n,r})
   &=\mathbf1_{\{r=9\}}u_{d_{n,r}+1},
 \label{eq:memoria-arista-ch}\\
       \operatorname{Carr}(e_{\varepsilon,n,r})
   &=d_{\varepsilon,r}\kappa_{\rm Carr},&
 d_{\varepsilon,r}
   &:=\mathbf1_{\{r=\mathfrak p_\varepsilon^{(2)}\}}
       \chi\!\left(m(\mathfrak p_\varepsilon^{(1)}),
                    m(\mathfrak p_\varepsilon^{(2)})\right).
 \label{eq:carry-arista-derivado-ch}
\end{align}
For identities and every admissible composition, also define
\begin{align}
 \mathsf C_{\mathsf M}(\operatorname{id}_{\widetilde q};d)
   &:=\operatorname{id}_{\mathsf M_d},&
 \kappa_{\mathsf M}(\operatorname{id}_{\widetilde q};d)&:=0,
 \label{eq:memoria-identidad-ch}\\
 \mathsf C_{\mathsf M}(e_t\cdots e_1)
   &:=\mathsf C_{\mathsf M}(e_t)\circ\cdots\circ
      \mathsf C_{\mathsf M}(e_1),&
 \kappa_{\mathsf M}(e_2e_1)
   &:=\mathsf C_{\mathsf M}(e_2)
        \bigl(\kappa_{\mathsf M}(e_1)\bigr)
      +\kappa_{\mathsf M}(e_2).
 \label{eq:cociclo-memoria-calibrado-ch}
\end{align}
The identity carries degree \(d\) because the same observable configuration
may recur at different depths. Iterated, the final equality defines
\(\kappa_{\mathsf M}(e_t\cdots e_1)\) for every length and is the cocycle law
of the direct system; the inclusions satisfy
\(\jmath_{d+2,d+1}\jmath_{d+1,d}=\jmath_{d+2,d}\).
In addition, set \(s(e_{\varepsilon,n,r})=\top\). The following twenty-seven
entries certify the coordinates distinguishing the three extensions; the
universal transport fields are specified immediately afterward. The symbols
\(I\) and \(II\) designate, respectively, the first retained emission and the
second emission performing the reduction; a dash denotes a zero carry
increment, not the absence of an edge.
\begin{center}
\small
\setlength{\tabcolsep}{3.2pt}
\begin{tabular}{c c cc cc cc c}
\toprule
\(r\)&\(M_{\rm ph}(r)\)&\multicolumn{2}{c}{\(\mathfrak p_{+1}\)}&
\multicolumn{2}{c}{\(\mathfrak p_0\)}&
\multicolumn{2}{c}{\(\mathfrak p_{-1}\)}&
\(\Delta\operatorname{mem}\)\\
& &mark&\(d_{+1,r}\)&mark&\(d_{0,r}\)&mark&\(d_{-1,r}\)&\\
\midrule
1&\(+1\)&I&0&--&0&--&0&0\\
2&\(-1\)&--&0&--&0&I&0&0\\
3&0&--&0&I&0&--&0&0\\
4&\(+1\)&II&\(+1\)&--&0&--&0&0\\
5&\(-1\)&--&0&--&0&II&\(-1\)&0\\
6&0&--&0&II&0&--&0&0\\
7&\(+1\)&--&0&--&0&--&0&0\\
8&\(-1\)&--&0&--&0&--&0&0\\
9&0&--&0&--&0&--&0&\(u\)\\
\bottomrule
\end{tabular}
\end{center}
Thus, \(d_{\varepsilon,r}\) is computed from two outputs of the selector and
is distinct from the cell quotients \(q^\Sigma,q^\Pi\). The ledger also
preserves the marked pair and the balanced residue; in particular, the zero
branch is not confused with an absence of event.

The carry composition law used by the TPK category is
\begin{equation}
 \operatorname{Carr}(\operatorname{id}_{\widetilde q})=0,\qquad
 \operatorname{Carr}(e_2e_1)
 =H_{\rm Carr}(e_1)\bigl(\operatorname{Carr}(e_2)\bigr)
  +Y_{\rm Carr}(e_2)\bigl(\operatorname{Carr}(e_1)\bigr).
 \label{eq:ley-carry-correcta-ch}
\end{equation}
On identities and paths, they are extended functorially:
\[
\begin{aligned}
 H_{\rm Carr}(\operatorname{id})
   &=Y_{\rm Carr}(\operatorname{id})=\operatorname{id},\\
 H_{\rm Carr}(e_t\cdots e_1)
   &=H_{\rm Carr}(e_1)\circ\cdots\circ H_{\rm Carr}(e_t),\\
 Y_{\rm Carr}(e_t\cdots e_1)
   &=Y_{\rm Carr}(e_t)\circ\cdots\circ Y_{\rm Carr}(e_1).
\end{aligned}
\]
In this model all these actions remain the identity, but their typing makes
explicit which endomorphism acts on each summand of
\eqref{eq:ley-carry-correcta-ch}. On the calibrated edges,
\(H_{\rm Carr}=Y_{\rm Carr}=\operatorname{id}\), so composition adds the
increments already produced. Equation~\eqref{eq:ley-carry-correcta-ch}, not an
identification between orientation and quotient, governs the accumulated
ledger.

\paragraph{Independent construction of the transition tuples.}
A fiber state is written, in the order of its relevant fields, as
\[
 x=(a;q;\tau,o,h;b_{\rm ref};
       R^\Sigma,q^\Sigma,R^\Pi,q^\Pi;c;
       \operatorname{Sig};C,\partial C;p_{\APP},p_{\TRIT};
       \lambda;\operatorname{gr}_{\mathsf M};\operatorname{mem};
       s;\operatorname{Led}).
\]
The update notation \(x[f\leftarrow v]\) changes only the indicated field.
For \(h\in\{(\Sigma,\Pi),(\Pi,\Sigma)\}\), define
\[
 h_r(h)=
 \begin{cases}
  (\Sigma,\Pi),&r\in\{1,4,7\},\\
  (\Pi,\Sigma),&r\in\{2,5,8\},\\
 h,&r\in\{3,6,9\}.
 \end{cases}
\]
The complete sheet-and-arithmetic coordinate is abbreviated as
\[
 \widehat h(x):=
 \bigl(h(x),R^\Sigma(x),q^\Sigma(x),R^\Pi(x),q^\Pi(x)\bigr).
\]
Let \(\widehat{\mathcal H}^{\rm cal}_{\widetilde q_{n,r}}\) be the set of
quintuples
\[
 \bigl(h,R^\Sigma_{a_{n,r}},q^\Sigma_{a_{n,r}},
          R^\Pi_{a_{n,r}},q^\Pi_{a_{n,r}}\bigr),
 \qquad
 h\in\{(\Sigma,\Pi),(\Pi,\Sigma)\},
\]
whose four arithmetic data are the two Euclidean divisions of the labeled APP
cell \(a_{n,r}\). This direct definition does not presuppose a state in a
carrier not yet constructed. Let
\(\mathsf{Hist}^{\rm cal}_{\widetilde q}\) be the set of admissible finite
paths in the graph of labeled edges ending at \(\widetilde q\), including the
identity, and let \(\mathsf B^{\rm cal}_{\widetilde q}\) be the set of their
ordered endpoint sequences. Define
\[
 \mathsf S^{\rm cal}_{\widetilde q}:=
 \widehat{\mathcal H}^{\rm cal}_{\widetilde q}
 \times\{o_{\rightarrow},o_{\circ},o_{\leftarrow}\}
 \times K_{\rm Carr}\times
 \mathsf{Hist}^{\rm cal}_{\widetilde q}\times
 \mathsf B^{\rm cal}_{\widetilde q}\times\mathsf M .
\]
The calibrated signature is not postulated: it is the typed assembly map
\[
 \operatorname{Sig}^{\rm cal}_{\widetilde q}:
 \widehat{\mathcal H}^{\rm cal}_{\widetilde q}
 \times\{o_{\rightarrow},o_{\circ},o_{\leftarrow}\}
 \times K_{\rm Carr}\times
 \mathsf{Hist}^{\rm cal}_{\widetilde q}\times
 \mathsf B^{\rm cal}_{\widetilde q}\times\mathsf M
 \longrightarrow\mathsf S^{\rm cal}_{\widetilde q},
\quad
 (\widehat h,o,c,\lambda,\partial C,\mu)
 \longmapsto(\widehat h,o,c,\lambda,\partial C,\mu).
\]
Likewise, \(\pi_{\mathcal F}\) denotes the projection that removes only the
first two coordinates \((a;q)\) of the state tuple; it preserves
\(\tau,o,h,b_{\rm ref}\), the four Euclidean data, carry, signature, cylinder,
boundary, prefixes, route, memory degree, memory, survival, and ledger. The
corresponding orientation is fixed by
\(o(+1)=o_{\rightarrow}\), \(o(0)=o_{\circ}\), and
\(o(-1)=o_{\leftarrow}\).
The calibrated APP fiber over \(\widetilde q_{n,r}\) contains the already
defined labeled copy \(a_{n,r}\). The initial state is fixed without free
coordinates:
\[
 \widehat h_\star:=
 \bigl((\Sigma,\Pi),R^\Sigma_{1,9},q^\Sigma_{1,9},
       R^\Pi_{1,9},q^\Pi_{1,9}\bigr).
\]
\begin{equation}
\begin{aligned}
 x_\star:=\bigl(&a_{0,1};\widetilde q_{0,1};
 +1,o_{\rightarrow},(\Sigma,\Pi);\bot;
 R^\Sigma_{1,9},q^\Sigma_{1,9},R^\Pi_{1,9},q^\Pi_{1,9};0;\\
 &\operatorname{Sig}^{\rm cal}_{\widetilde q_{0,1}}
       (\widehat h_\star,o_{\rightarrow},0,
        \operatorname{id}_{\widetilde q_{0,1}},\varnothing,
        \jmath_{k_0}(0_{\mathsf M_{k_0}}));
 \varnothing,\varnothing;
 (a_{0,1}),(+1);\operatorname{id}_{\widetilde q_{0,1}};
 k_0;0_{\mathsf M_{k_0}};\top;\varnothing\bigr).
\end{aligned}
\label{eq:estado-inicial-calibrado-ch}
\end{equation}
The tuple contains no indeterminate coordinates and satisfies
\(s(x_\star)=\top\). Let \(G_0:=\{x_\star\}\). Assuming \(G_n\) has been
constructed, fix \(x\in G_n\), \(\varepsilon\in\{-1,0,+1\}\), and
\(x_{\varepsilon,0}(x):=x\). For \(r=1,\ldots,9\), first construct the
elementary ledger
\begin{equation}
\begin{aligned}
 \ell_{\varepsilon,n,r}:=\bigl(&e_{\varepsilon,n,r},a_{n,r},a_{(n,r)^+},
 \widetilde q_{n,r},\widetilde q_{n,r+1},M_{\rm ph}(r),\\
 &h_r(h(x_{\varepsilon,r-1}(x))),o(M_{\rm ph}(r)),\mathfrak p_\varepsilon,
 d_{\varepsilon,r},\\
 &\mathsf C_{\mathsf M}(e_{\varepsilon,n,r}),
 \kappa_{\mathsf M}(e_{\varepsilon,n,r}),
 \operatorname{Carr}(e_{\varepsilon,n,r})\bigr),
\end{aligned}
 \label{eq:registro-elemental-ext-ch}
\end{equation}
and set
\(\operatorname{Led}^{\rm new}_{\varepsilon,n,r}:=
\operatorname{Led}(x_{\varepsilon,r-1}(x))^\frown
\ell_{\varepsilon,n,r}\). Then construct the tuple
\begin{equation}
 y_{\varepsilon,r}(x):=
 x_{\varepsilon,r-1}(x)
 [\tau\leftarrow M_{\rm ph}(r),\
  h\leftarrow h_r(h(x_{\varepsilon,r-1}(x))),\
  o\leftarrow o(M_{\rm ph}(r)),\
  b_{\rm ref}\leftarrow\mathfrak p_\varepsilon],
 \label{eq:tupla-sel-cal-ch}
\end{equation}
where selection leaves the cell, quotients, carry, route, memory, cylinder,
boundary, survival, and ledger unchanged; the coordinate \(b_{\rm ref}\)
makes the chosen balanced reduction explicit. Membership in an ambient
relation is not presupposed here.

Next, define \(z_{\varepsilon,r}(x)\) field by field: its configuration is
\(\widetilde q_{n,r+1}\); its complete cell is \(a_{(n,r)^+}\), with the four
residues and quotients recomputed by the two Euclidean divisions; it preserves
the selected sheet and orientation; and it performs the following six updates,
in which only the common argument \(x\) is omitted:
\begin{equation}
\begin{aligned}
 C(z_{\varepsilon,r})&:=C(y_{\varepsilon,r})^\frown
       e_{\varepsilon,n,r},\\
 \partial C(z_{\varepsilon,r})&:=\partial C(y_{\varepsilon,r})^\frown
       (\widetilde q_{n,r},\widetilde q_{n,r+1}),\\
 p_{\APP}(z_{\varepsilon,r})&:=p_{\APP}(y_{\varepsilon,r})^\frown
       a_{(n,r)^+},\\
 p_{\TRIT}(z_{\varepsilon,r})&:=p_{\TRIT}(y_{\varepsilon,r})^\frown
       \tau(r),\\
 \operatorname{Led}(z_{\varepsilon,r})&:=
       \operatorname{Led}^{\rm new}_{\varepsilon,n,r},\\
 \lambda(z_{\varepsilon,r})&:=e_{\varepsilon,n,r}
       \lambda(y_{\varepsilon,r}),
\end{aligned}
 \label{eq:tupla-tra-cal-ch}
\end{equation}
and still retains the fields
\(c,\operatorname{gr}_{\mathsf M},\operatorname{mem},s\) of
\(y_{\varepsilon,r}(x)\). Its signature, by contrast, is retyped at the target
configuration through the transported presignature
\begin{equation}
\begin{aligned}
 \operatorname{Sig}(z_{\varepsilon,r}(x))
  :=\operatorname{Sig}^{\rm cal}_{\widetilde q_{n,r+1}}
       \Bigl(&\widehat h(z_{\varepsilon,r}(x)),
             o(z_{\varepsilon,r}(x)),c(y_{\varepsilon,r}(x)),
             \lambda(z_{\varepsilon,r}(x)),\\[-2pt]
             &\partial C(z_{\varepsilon,r}(x)),
             \jmath_{d_{n,r}}
                (\operatorname{mem}_{d_{n,r}}
                   (y_{\varepsilon,r}(x)))\Bigr).
\end{aligned}
\label{eq:prefirma-transporte-cal-ch}
\end{equation}
Thus, and only in \(\mathsf{Tra}^{\rm cal}\), the edge, its endpoints, the
mark \((\mathfrak p_\varepsilon,r,d_{\varepsilon,r})\), and its cocycles are
added to the incidence ledger.

The initial state satisfies
\[
 c(x_\star)=0=\operatorname{Carr}
   (\operatorname{id}_{\widetilde q_{0,1}})
   =\operatorname{Carr}(\lambda(x_\star)).
\]
In the complete states \(x_{\varepsilon,r-1}(x)\), and also after the
selection \(y_{\varepsilon,r}(x)\), the typed invariant
\(c(v)=\operatorname{Carr}(\lambda(v))\) is preserved, since
\(\mathsf{Sel}^{\rm cal}\) modifies neither field. Transport advances the
route but leaves the carry pending; hence the pre-update tuple satisfies
exactly
\[
 c(z_{\varepsilon,r}(x))=c(y_{\varepsilon,r}(x))
 =\operatorname{Carr}(\lambda(y_{\varepsilon,r}(x))).
\]
The map \(\operatorname{Upd}^{\rm cal}\) then restores
\(c(x_{\varepsilon,r}(x))=
\operatorname{Carr}(\lambda(x_{\varepsilon,r}(x)))\) by means of the
composition law, without attributing that invariant to the intermediate state \(z\).

Finally, apply the unary function
\(\operatorname{Upd}^{\rm cal}_{e_{\varepsilon,n,r}}
=\mathsf{Mem}^{\rm cal}_{e_{\varepsilon,n,r}}\) to that tuple. Its result
\(x_{\varepsilon,r}(x)\) is given by
\begin{equation}
\begin{aligned}
 c\bigl(x_{\varepsilon,r}(x)\bigr)
   &=\operatorname{Carr}\bigl(\lambda(z_{\varepsilon,r}(x))\bigr)\\
   &=H_{\rm Carr}(\lambda(y_{\varepsilon,r}(x)))
       \bigl(\operatorname{Carr}(e_{\varepsilon,n,r})\bigr)
     +Y_{\rm Carr}(e_{\varepsilon,n,r})
       \bigl(c(y_{\varepsilon,r}(x))\bigr),\\
 \lambda\bigl(x_{\varepsilon,r}(x)\bigr)
   &=\lambda(z_{\varepsilon,r}(x)),\\
 \operatorname{gr}_{\mathsf M}
      \bigl(x_{\varepsilon,r}(x)\bigr)
   &=d_{n,r}+1,\\
 \operatorname{mem}_{d_{n,r}+1}
       \bigl(x_{\varepsilon,r}(x)\bigr)
   &=\mathsf C_{\mathsf M}(e_{\varepsilon,n,r})
       \operatorname{mem}_{d_{n,r}}(z_{\varepsilon,r}(x))
     +\kappa_{\mathsf M}(e_{\varepsilon,n,r}),\\
 s\bigl(x_{\varepsilon,r}(x)\bigr)
   &=s(e_{\varepsilon,n,r})\wedge s(z_{\varepsilon,r}(x)),\\
 \operatorname{Sig}\bigl(x_{\varepsilon,r}(x)\bigr)
   &=\operatorname{Sig}^{\rm cal}_{\widetilde q_{n,r+1}}
       \Bigl(\widehat h(z_{\varepsilon,r}(x)),
             o(z_{\varepsilon,r}(x)),
             c(x_{\varepsilon,r}(x)),
             \lambda(z_{\varepsilon,r}(x)),
             \partial C(z_{\varepsilon,r}(x)),\\[-2pt]
   &\hspace{112pt}
             \jmath_{d_{n,r}+1}
              \bigl(\operatorname{mem}_{d_{n,r}+1}
                 (x_{\varepsilon,r}(x))\bigr)\Bigr).
\end{aligned}
\label{eq:upd-unaria-tres-vueltas-ch}
\end{equation}
All fields not enumerated in \eqref{eq:upd-unaria-tres-vueltas-ch} coincide
with those of \(z_{\varepsilon,r}(x)\). In particular, the signature is
computed from the sheet, orientation, boundary, and route already fixed by
\(\mathsf{Sel}^{\rm cal}\) and \(\mathsf{Tra}^{\rm cal}\), and from the carry
and memory computed in the preceding lines; it is not circularly defined from
the final state. Neither \(\varepsilon\), nor \(d_{\varepsilon,r}\), nor
\(u\) is an additional argument of \(\operatorname{Upd}^{\rm cal}\): the
function reads all three data from the already constructed edge.

\paragraph{Simultaneous induction of the levels.}
Once the preceding nine stages have been completed for an already constructed
\(G_n\), define immediately
\begin{equation}
 \gamma_{\varepsilon,k_n}(x):=x_{\varepsilon,9}(x),\qquad
 G_{n+1}:=\{\gamma_{\varepsilon,k_n}(x):
       x\in G_n,\ \varepsilon\in\TRIT\}.
 \label{eq:injerto-total-tres-vueltas-ch}
\end{equation}
\label{eq:tres-vueltas-tpk-ch}
The base \(G_0=\{x_\star\}\), the intermediate tuples, and
\eqref{eq:injerto-total-tres-vueltas-ch} constitute a single definition by
recursion on \(n\). Therefore, no subsequent carrier is used to assume the
existence of \(G_n\): the next level is introduced only after constructing all
its edges from the preceding level.

\paragraph{TPK extension relation and effective membership.}
Let \(\mathcal E^{\rm cal}_{\widetilde q}\) be the least labeled collection
containing \(x_\star\) and, at each step of the preceding recursion, the four
tuples
\(x_{\varepsilon,r-1}(x),y_{\varepsilon,r}(x),
z_{\varepsilon,r}(x),x_{\varepsilon,r}(x)\) in their source and target
configurations. Also let
\(\mathcal A^{\rm cal}_{\widetilde q_{n,r}}:=\{a_{n,r}\}\), and set
\[
 \widetilde{\mathcal Q}_{\rm obs}:=
   \{\widetilde q_{n,r}:n\ge0,\ 1\le r\le9\},
 \qquad
 \mathcal E^{\rm cal}:=\bigsqcup_{\widetilde q}
     \mathcal E^{\rm cal}_{\widetilde q},
 \qquad
 \mathcal F^{\rm cal}_{\widetilde q}:=
     \pi_{\mathcal F}(\mathcal E^{\rm cal}_{\widetilde q}).
\]
For each depth \(d\), let
\(\mathcal E^{\rm cal}_d:=
\operatorname{gr}_{\mathsf M}^{-1}(d)\subset\mathcal E^{\rm cal}\). The
explicit degree coordinate, rather than the mere membership of a
representative in some image of the direct system, makes these carriers
disjoint. On them, the projections are typed as
\[
\begin{aligned}
 \operatorname{cfg}&:\mathcal E^{\rm cal}\longrightarrow
     \widetilde{\mathcal Q}_{\rm obs},\\
 g(x)&:=\operatorname{fase}\!
       \bigl(\pi_{\rm obs}(\operatorname{cfg}(x))\bigr)\in C_9,\\
 c&:\mathcal E^{\rm cal}\longrightarrow K_{\rm Carr},\\
 s&:\mathcal E^{\rm cal}\longrightarrow\{\bot,\top\},\\
 \lambda&:\mathcal E^{\rm cal}_{\widetilde q}
       \longrightarrow\mathsf{Hist}^{\rm cal}_{\widetilde q},\\
 \operatorname{gr}_{\mathsf M}&:\mathcal E^{\rm cal}
       \longrightarrow\mathbb N_{\ge1},\\
 \operatorname{mem}_d&:\mathcal E^{\rm cal}_d
       \longrightarrow\mathsf M_d,\\
 \operatorname{mem}(x)&:=
       \jmath_d\bigl(\operatorname{mem}_d(x)\bigr)\in\mathsf M
       \quad(x\text{ of depth }d).
\end{aligned}
\]
Here \(\operatorname{cfg}(x)\) is the second coordinate of the tuple, and
\(\operatorname{fase}\) reads the phase of the observable calendar. For the
three carrier classes, fix the diagonals
\[
\begin{aligned}
 \Delta_{\APP,\widetilde q}
   &:=\{(a,a):a\in\mathcal A^{\rm cal}_{\widetilde q}\},\\
 \Delta_{{\rm fib},\widetilde q}
   &:=\{(f,f):f\in\mathcal F^{\rm cal}_{\widetilde q}\},\\
 \Delta_{{\rm enr},\widetilde q}
   &:=\{(x,x):x\in\mathcal E^{\rm cal}_{\widetilde q}\}.
\end{aligned}
\]
In particular,
\[
 a_{0,1}\in\mathcal A^{\rm cal}_{\widetilde q_{0,1}},
 \qquad
 x_\star\in\mathcal E^{\rm cal}_{\widetilde q_{0,1}}.
\]
On these carriers, the following relations are defined by their complete
graphs, and not through necessary implications:
\begin{align}
 \mathsf{Sel}^{\rm cal}_{e_{\varepsilon,n,r}}
   &:=\{(x_{\varepsilon,r-1}(x),y_{\varepsilon,r}(x)):
          x\in G_n\},\\
 \mathsf{Tra}^{\rm cal}_{e_{\varepsilon,n,r}}
   &:=\{(y_{\varepsilon,r}(x),z_{\varepsilon,r}(x)):
          x\in G_n\},\\
 \operatorname{Ext}^{\APP,\rm cal}_{e_{\varepsilon,n,r}}
   &:=\{(a_{n,r},a_{(n,r)^+})\},
 \label{eq:ext-app-calibrada-ch}\\
 \operatorname{Ext}^{\rm fib,cal}_{e_{\varepsilon,n,r}}
   &:=\{(\pi_{\mathcal F}x_{\varepsilon,r-1}(x),
          \pi_{\mathcal F}x_{\varepsilon,r}(x)):x\in G_n\}.
 \label{eq:ext-fibra-calibrada-ch}
\end{align}
The synchronized enriched extension is
\begin{equation}
 \operatorname{Ext}^{\rm enr,cal}_{e_{\varepsilon,n,r}}
 :=\{(x_{\varepsilon,r-1}(x),x_{\varepsilon,r}(x)):
       x\in G_n\}.
 \label{eq:ext-enriquecida-calibrada-ch}
\end{equation}
The third map is defined by
\begin{equation}
 \operatorname{graph}
 \bigl(\operatorname{Upd}^{\rm cal}_{e_{\varepsilon,n,r}}\bigr)
 :=\{(z_{\varepsilon,r}(x),x_{\varepsilon,r}(x)):
       x\in G_n\},
 \label{eq:grafo-upd-calibrado-ch}
\end{equation}
where the second component is exactly the one in
\eqref{eq:upd-unaria-tres-vueltas-ch}. Consequently,
\begin{equation}
 \mathcal U^{(1),\rm cal}_{e_{\varepsilon,n,r}}
 :=\operatorname{Upd}^{\rm cal}_{e_{\varepsilon,n,r}}\circ
   \mathsf{Tra}^{\rm cal}_{e_{\varepsilon,n,r}}\circ
   \mathsf{Sel}^{\rm cal}_{e_{\varepsilon,n,r}}
 =\{(x_{\varepsilon,r-1}(x),x_{\varepsilon,r}(x)):
      x\in G_n\}.
 \label{eq:transicion-cal-en-tpk-ch}
\end{equation}
Each operator is indexed by the edge \(e_{\varepsilon,n,r}\) produced by the
reduction \(\mathfrak p_\varepsilon\), rather than by a trit supplied from
outside. The coordinate \(b_{\rm ref}\) also distinguishes their codomains.
Therefore,
\(\mathsf{Sel}^{\rm cal}_{e}\),
\(\mathsf{Tra}^{\rm cal}\) and \(\operatorname{Upd}^{\rm cal}\) are functions
on their respective domains, whereas the union
\(\bigcup_{\varepsilon\in\TRIT}
\mathcal U^{(1),\rm cal}_{e_{\varepsilon,n,r}}\) deliberately preserves the
three outputs of the trisection.

For \(\bullet\in\{\APP,{\rm fib},{\rm enr}\}\), identities and compositions
are not left implicit. Define
\begin{align}
 \operatorname{Ext}^{\bullet,\rm cal}_{\operatorname{id}_{\widetilde q}}
   &:=\Delta_{\bullet,\widetilde q},\\
 \operatorname{Ext}^{\bullet,\rm cal}_{e_t\cdots e_1}
   &:=\operatorname{Ext}^{\bullet,\rm cal}_{e_t}\circ\cdots\circ
      \operatorname{Ext}^{\bullet,\rm cal}_{e_1}
 \label{eq:ext-palabra-calibrada-ch}
\end{align}
for every admissible edge word. The two APP and fiber path closures are formed
separately,
\begin{equation}
 \operatorname{Path}(\operatorname{Ext}^{\APP,\rm cal}),\qquad
 \operatorname{Path}(\operatorname{Ext}^{\rm fib,cal}),
 \label{eq:dos-clausuras-ext-ch}
\end{equation}
and \(\mathsf{Tra}^{\rm cal}\) synchronizes them edge by edge.

Finally, the levels of complete states and their truncations are fixed by the
recursion itself:
\begin{align}
 H^{\rm cal}_{d_{n,r}}
   &:=\{x_{\varepsilon,r-1}(x):
          x\in G_n,\ \varepsilon\in\TRIT\},\\
 H^{\rm cal}_{d_{n,r}+1}
   &:=\{x_{\varepsilon,r}(x):
          x\in G_n,\ \varepsilon\in\TRIT\},\\
 \rho^{\rm cal}_{d_{n,r}+1,d_{n,r}}
   \bigl(x_{\varepsilon,r}(x)\bigr)
   &:=x_{\varepsilon,r-1}(x).
 \label{eq:truncamiento-calibrado-definido-ch}
\end{align}
The final ledger preserves \((n,r,\mathfrak p_\varepsilon)\), so the
predecessor of the left-hand side is unique and \(\rho^{\rm cal}\) is
well-defined. For \(b>a\), set
\(\rho^{\rm cal}_{b,a}:=\rho^{\rm cal}_{a+1,a}\circ\cdots\circ
\rho^{\rm cal}_{b,b-1}\). This map recovers the complete preceding state,
including its signature; it does not purport to invert an overwritten
coordinate without consulting the ledger.

\begin{proposition}[Autonomous realization of the TPK through explicit transitions]
\label{prop:modelo-tpk-calibrado-ch}
The collection
\[
\begin{aligned}
 \mathfrak T_{\rm TPK}^{\rm cal}:=\bigl(
 &(\widetilde q_{n,r})_{n,r},\pi_{\rm obs},
   \widetilde F_{\rm obs},M_{\rm ph},\\
 &(\mathcal A_{\widetilde q}^{\rm cal})_{\widetilde q},
   (\mathcal E_{\widetilde q}^{\rm cal})_{\widetilde q},
   (\mathcal F_{\widetilde q}^{\rm cal})_{\widetilde q},
   \operatorname{Ext}^{\APP,\rm cal},
   \operatorname{Ext}^{\rm fib,cal},
   \operatorname{Ext}^{\rm enr,cal},\rho^{\rm cal},\\
 &\mathsf{Sel}^{\rm cal},\mathsf{Tra}^{\rm cal},
   \operatorname{Upd}^{\rm cal},\mathsf C_{\mathsf M},
   \kappa_{\mathsf M},\operatorname{gr}_{\mathsf M},
   H_{\rm Carr},Y_{\rm Carr},\\
 &\operatorname{Sig}^{\rm cal},\pi_{\mathcal F}\bigr).
\end{aligned}
\]
is an enriched model of the TPK. In particular, every pair on the right-hand
side of \eqref{eq:transicion-cal-en-tpk-ch} is an effective TPK edge, not a
consequence inferred from a merely necessary condition.
\end{proposition}

\begin{proof}
The labels \((n,r)\) make the carrier copies disjoint and unambiguously fix
each source and target. The relation \(\operatorname{Ext}^{\APP,\rm cal}\)
transports the complete cell and preserves the two Euclidean identities; the
relation \(\operatorname{Ext}^{\rm fib,cal}\) relates the fibers of the
preceding and subsequent complete states: it transports sheet, orientation,
prefix, cylinder, boundary, and ledger, and incorporates the updated carry,
memory, and signature. Selection does not modify the ledger, and transport
concatenates exactly one entry and produces the target presignature
\eqref{eq:prefirma-transporte-cal-ch}. The update recomputes the signature by
means of the typed map \(\operatorname{Sig}^{\rm cal}_{\widetilde q}\) from
the already determined sheet, orientation, carry, route, boundary, and memory.
The map \(\operatorname{Upd}^{\rm cal}\) is unary because the edge contained
in \(z_{\varepsilon,r}(x)\) determines its cocycles.

The identities are the diagonals, and the word extensions are the compositions
in \eqref{eq:ext-palabra-calibrada-ch}. Concatenation of prefixes and ledgers
is associative; memory satisfies the cocycle law, and carry satisfies
\eqref{eq:ley-carry-correcta-ch}. Therefore, the two parenthesizations of three
edges produce the same state. Formula
\eqref{eq:truncamiento-calibrado-definido-ch} defines truncation on complete
states, and uniqueness of the final ledger proves that it recovers the source
in every field. These are the identity, source, target, composition, memory,
carry, and truncation laws of the enriched TPK. In particular, the extension
relations are determined as graphs of effective maps rather than as
consequences of merely necessary conditions.
\end{proof}

Each \(\gamma_{\varepsilon,k_n}\) comprises exactly nine applications
of \(\operatorname{Upd}^{\rm cal}_{e_{\varepsilon,n,r}}\circ
\mathsf{Tra}^{\rm cal}_{e_{\varepsilon,n,r}}\circ
\mathsf{Sel}^{\rm cal}_{e_{\varepsilon,n,r}}\). Survival will be proved by
means of a constructed tail and will not be used to define \(G_n\).

\begin{lemma}[Effective internal trisection of one turn]
\label{lem:elevacion-catalogal-nonadica}
For every \(n\ge0\), every \(x\in G_n\), and every
\(\varepsilon\in\{-1,0,+1\}\), the expression
\(\gamma_{\varepsilon,k_n}(x)\) is defined by nine TPK transitions, with its
synchronized relations
\(\operatorname{Ext}^{\APP,\rm cal}\) and
\(\operatorname{Ext}^{\rm fib,cal}\), and satisfies
\begin{align}
 \rho^{\rm cal}_{k_n+9,k_n}\gamma_{\varepsilon,k_n}(x)&=x,
 \label{eq:truncamiento-injerto-ch}\\
 g\bigl(\gamma_{\varepsilon,k_n}(x)\bigr)
   &=g(x),&
 \jmath_{k_n+9}\!\left(\operatorname{mem}_{k_n+9}
        (\gamma_{\varepsilon,k_n}(x))\right)
   &=\jmath_{k_n}\!\left(\operatorname{mem}_{k_n}(x)\right)+u,
 \label{eq:fase-memoria-injerto-ch}\\
 c\bigl(\gamma_{\varepsilon,k_n}(x)\bigr)
   &=c(x)+\varepsilon\kappa_{\rm Carr}.&&
 \label{eq:carry-injerto-ch}
\end{align}
\label{eq:tres-vueltas-propiedades-ch}
\label{eq:tres-vueltas-acarreo-ch}
The ordered ledger of the nine edges determines \(\varepsilon\) uniquely.
The three images are therefore disjoint. Each endpoint belongs to
\(G_{n+1}\) and again admits the three maps
\(\gamma_{-1,k_{n+1}},\gamma_{0,k_{n+1}},
\gamma_{+1,k_{n+1}}\).
\end{lemma}

\begin{proof}
Formulas \eqref{eq:celula-completa-fase-ch} traverse the nine phases once and
return the copy \(a_{n+1,1}\) of \(a(1)\). At each step, the APP relation
transports the complete sextuple; \(\mathsf{Sel}^{\rm cal}\) fixes the TRIT
type from the phase mark \(r\) preserved in the state; and
\(\mathsf{Tra}^{\rm cal}\) adds exactly one symbol to the prefix, one segment
to the cylinder, and its ordered pair of endpoints to the boundary. By
induction on \(r\), the explicit tuples \(y_{\varepsilon,r}(x)\) and
\(z_{\varepsilon,r}(x)\) satisfy the predicates of the indicated relations,
and the unary update produces the state at the next level. The identity
\(\rho^{\rm cal}_{d_{n,r}+1,d_{n,r}}
(x_{\varepsilon,r}(x))=x_{\varepsilon,r-1}(x)\) is the definition
\eqref{eq:truncamiento-calibrado-definido-ch}; its well-definedness follows
from uniqueness of the final ledger. The composition of the nine truncation
maps recovers \(x\), proving \eqref{eq:truncamiento-injerto-ch} without
separately inverting the overwritten fields.

The phase \(9\to1\) occurs exactly once. By
\eqref{eq:memoria-arista-ch} and the cocycle law, the linear memory component
transports the preceding memory, and its translational component adds exactly
\(u\) in the direct limit. The compatibility
\(\jmath_{d+1}\jmath_{d+1,d}=\jmath_d\) proves
\eqref{eq:fase-memoria-injerto-ch}; the return concerns the phase, not the
enriched state.

By \eqref{eq:carry-arista-derivado-ch}, eight edges have zero increment, and
the edge completing \(\mathfrak p_\varepsilon\) has the computed increment
\(\chi(\varepsilon,\varepsilon)\kappa_{\rm Carr}\). The laws
\eqref{eq:ley-carry-correcta-ch} and
\eqref{eq:carry-no-parametrico-ch}
give \eqref{eq:carry-injerto-ch}. The ledger preserves
\((\mathfrak p_\varepsilon,2\varepsilon,r_3(2\varepsilon),
\chi(\varepsilon,\varepsilon))\); the three values of this ledger are
distinct even when the scalar increment is zero. Hence the images are not
identified, and the block decoder recovers a unique \(\varepsilon\).

The edge definitions depend only on the new phase and concatenate, rather than
replace, the prefix, cylinder, boundary, and ledger. They are also invariant
under
\(\jmath_d\operatorname{mem}_d\mapsto
  \jmath_d\operatorname{mem}_d+nu\)
and under translation of the input carry. Consequently, the same construction
applies at the endpoint of any turn. For every finite state, the explicit
sequence
\[
 x,\quad\gamma_{0,k_n}(x),\quad
 \gamma_{0,k_{n+1}}\gamma_{0,k_n}(x),\quad\ldots
\]
is a compatible infinite tail. This proves the survival of all constructed
states and the repeated availability of the three branches, without
postulating either property. Consequently, in accordance with the preceding
level definitions, we may finally identify
\begin{equation}
 H_{k_n}^{\rm cal}:=G_n
 \label{eq:subarbol-calibrado-superviviente-ch}
\end{equation}
for every \(n\ge0\).
\end{proof}

The preceding closure distinguishes the three clocks involved. Nine emissions
return the phase in \(C_9\) and advance memory. Twelve turns, that is,
\(108\) emissions, also return the position of the observable calendar; even
then, neither the ledger nor memory is erased. The operator
\(\mathcal K_{\rm ph}\) acts only afterward, as a recognition of reduced
memory from the already constructed turns; it never generates the edges of
\(\operatorname{Ext}^{\rm enr,cal}\).


Let
\[
 \upsilon:\mathbb F_3\longrightarrow\{-1,0,+1\},\qquad
 \upsilon(0)=0,\quad\upsilon(1)=+1,\quad\upsilon(2)=-1
\]
be the balanced representative. For
\(b=(b_1,\ldots,b_6)\in\mathbb F_3^6\), the block extension is the
ordered composition of six turns
\[
 \Gamma_{b,k}^{[54]}:=
 \gamma_{\upsilon(b_6),k+45}\circ\gamma_{\upsilon(b_5),k+36}\circ
 \gamma_{\upsilon(b_4),k+27}\circ\gamma_{\upsilon(b_3),k+18}\circ
 \gamma_{\upsilon(b_2),k+9}\circ\gamma_{\upsilon(b_1),k}.
\tag{C.39d}
\label{eq:gamma-bloque-54-ch}
\]
At block level one always uses
\(k=k_{6N}=1+54N\); consequently, each factor has the domain constructed by
the preceding factor.
Fix a canonical APP seed \(x_\star\in H^{\rm cal}_1\) and define recursively
\[
 \begin{aligned}
  Y_0&:=\{x_\star\},\qquad Y_N\subseteq H_{1+54N}^{\rm cal},\\
  Y_{N+1}&:=\{\Gamma_{b,1+54N}^{[54]}(x):
                 x\in Y_N,\ b\in\mathbb F_3^6\},\\
  \rho^{\rm blk}_{N+1,N}
    &:=\rho^{\rm cal}_{1+54(N+1),1+54N}
       \!\upharpoonright Y_{N+1},\\
  \operatorname{Ext}^{\rm cal}_N(x)
    &:=\{\Gamma_{b,1+54N}^{[54]}(x):b\in\mathbb F_3^6\}.
 \end{aligned}
\tag{C.39e}
\label{eq:sincronizacion-54-ch}
\]
The notation \(\rho^{\rm cal}_{b,a}\) denotes the composition of all deletion maps
between levels \(b\) and \(a\). By construction, every extension in
\(\operatorname{Ext}^{\rm cal}_N(x)\) is a chain of fifty-four
TPK edges, not a jump of the Markov quotient.

The reader is defined directly on the ledger, before invoking a
representation through \(\Gamma_b\). The preceding lemma provides a
decoder
\[
 \operatorname{dec}_9:\{\text{generated nine-edge segments}\}
 \longrightarrow\{-1,0,+1\},
\]
which reads the TRIT pair and its reduction on the segment. If
\(\operatorname{Led}^{(9)}_{j,i}(x)\) denotes the \(i\)-th nonadic segment
of block \(j\) in the ledger of \(x\), set
\begin{equation}
 \beta_{{\rm blk},N}(x):=
 \Bigl(
   \bigl(\upsilon^{-1}\operatorname{dec}_9
          (\operatorname{Led}^{(9)}_{j,i}(x))\bigr)_{i=1}^{6}
 \Bigr)_{j=0}^{N-1}.
 \label{eq:lector-bloques-ch}
\end{equation}
For \(N=0\), the value is the empty word. This definition does not presuppose
that the decomposition of \(x\) as a composition of maps \(\Gamma_b\) is
unique.

We now specialize the preceding arithmetic conjugacy to
\(p=3\), \(d=6\), and \(A=I_6\). If
\(\pi_{N+1,N}:(\mathbb Z/3^{N+1}\mathbb Z)^6\to
(\mathbb Z/3^N\mathbb Z)^6\) is the canonical reduction, define
\begin{equation}
 \Phi_N:=\Psi_N^{-1}\circ\beta_{{\rm blk},N}:
 Y_N\longrightarrow(\mathbb Z/3^N\mathbb Z)^6.
 \label{eq:phi-hensel-bloques-ch}
\end{equation}
This definition does not assign a coordinate after observing a branch: the
reader \(\beta_{{\rm blk},N}\) is obtained from the already constructed TPK ledger, and
\(\Psi_N^{-1}\) is the explicit inverse, fixed beforehand, of the
unfiltered arithmetic tower.

\begin{proposition}[Finite conjugacy and complete covering of the fiber]
\label{prop:cobertura-729-seccion}
For every \(N\geq0\),
\[
 \beta_{{\rm blk},N}:Y_N\xrightarrow{\ \cong\ }
 (\mathbb F_3^6)^N
\]
is a bijection compatible with truncations, and for every \(x\in Y_N\)
\[
 \begin{gathered}
 \beta_{{\rm blk},N+1}
 \bigl[\operatorname{Ext}^{\rm cal}_N(x)\bigr]
 =\{\beta_{{\rm blk},N}(x)^\frown b:
       b\in\mathbb F_3^6\},\\
 \operatorname{pref}_N\circ\beta_{{\rm blk},N+1}
 =\beta_{{\rm blk},N}\circ\rho^{\rm blk}_{N+1,N}.
 \end{gathered}
\tag{C.40$_0$}
\label{eq:cobertura-729-ch}
\]
The maps \(\Phi_N\) are bijections and satisfy the intertwining identity
\begin{equation}
 \pi_{N+1,N}\circ\Phi_{N+1}
 =\Phi_N\circ\rho^{\rm blk}_{N+1,N}.
 \label{eq:entrelazamiento-hensel-ext-ch}
\end{equation}
Consequently, the tower of generated histories and the Hensel arithmetic
tower are isomorphic as finite inverse systems.
\end{proposition}

\begin{proof}
The assertion is immediate for \(N=0\). Assume that it holds at level
\(N\). For each word \(b\in\mathbb F_3^6\), the preceding lemma constructs an extension
    \(\Gamma_{b,1+54N}^{[54]}(x)\). If \(b\ne b'\), the first coordinate at
which they differ determines a distinct TRIT pair and cocycle in the first
divergent nonadic segment of the genealogical ledger. That ledger preserves
the ordered edges, not merely the final sum of their carries; hence the two
histories remain distinct even if carries from later turns cancel.
The direct reading \eqref{eq:lector-bloques-ch} then gives
\begin{equation}
 \beta_{{\rm blk},N+1}
 \bigl(\Gamma_{b,1+54N}^{[54]}(x)\bigr)
 =\beta_{{\rm blk},N}(x)^\frown b.
 \label{eq:beta-concatenacion-demostrada-ch}
\end{equation}
The six successive trisections produce exactly \(3^6\) distinct extensions.
Deleting the last fifty-four steps recovers \(x\), and deleting the final
block of the ledger recovers \(\beta_{{\rm blk},N}(x)\). This simultaneously
proves bijectivity and the compatibility of \(\beta_{{\rm blk},N}\) with
truncations.

Proposition~\ref{prop:ch-hensel-bruto} independently proves that
\(\Psi_N\) is bijective and that
\[
 \operatorname{pref}_N\circ\Psi_{N+1}
 =\Psi_N\circ\pi_{N+1,N}.
\]
Composing this equality with the compatibility of \(\beta_{\rm blk}\) and
applying \(\Psi_N^{-1}\) yields
\[
 \pi_{N+1,N}\Psi_{N+1}^{-1}\beta_{{\rm blk},N+1}
 =\Psi_N^{-1}\beta_{{\rm blk},N}\rho^{\rm blk}_{N+1,N},
\]
which is exactly \eqref{eq:entrelazamiento-hensel-ext-ch}. The equality does
not arise by defining an ad hoc coordinate on each branch: it results from
composing two systems of constructed bijections that are compatible at every
depth.
\end{proof}

The inverse limit
\[
 \Omega_\Gamma^{\rm cal}:=
 \varprojlim_N(Y_N,\rho^{\rm blk}_{N+1,N})
\tag{C.40a$_0$}
\label{eq:omega-gamma-catalogal-ch}
\]
is the history limit of the TPK model constructed edge by edge in the
preceding lemma. Its projection under \(\pi_{\rm obs}\) preserves the
configurations of the observable calendar, but the assertion does not depend
on identifying this autonomous model with an ambient relation defined only
by necessary conditions. The finite bijections pass to the limit and produce
the cylinder homeomorphism
\[
 \Lambda:(\mathbb F_3^6)^{\mathbb N}
 \xrightarrow{\ \cong\ }\Omega_\Gamma^{\rm cal},
 \qquad
 \beta_{\rm blk}:=\Lambda^{-1}.
\tag{C.40b$_0$}
\label{eq:homeomorfismo-calibrado-ch}
\]
Denote the canonical projection by
\(\rho^{\rm blk}_{\infty,N}:\Omega_\Gamma^{\rm cal}\to Y_N\).
Thus, \(\Lambda\) establishes an equivalence compatible with truncations:
each ternary tape encodes a TPK history, and every prefix is extended through
edges of the same generator.

% BEGIN REV09_RECTAS27_PROLONGACION
\subsubsection{Reversible prolongation in line coordinates}
\label{corpus9:xii:rectas-prolongacion}

Proposition~\ref{prop:cobertura-729-seccion} constructs the bijections
\(\beta_{{\rm blk},N}\) from nonadic segments of the ordered record,
together with their compatibility with block deletion. The homeomorphism
\eqref{eq:homeomorfismo-calibrado-ch} identifies the limit of these
generated histories with the sequence of their six-trit words.
We now compose that reader with the finite chart
\eqref{corpus9:xii:rectas-codificacion}, retaining its central law.

\begin{theorem}[Cylindrical transport of histories to ordered pairs of lines]
\label{corpus9:xii:rectas-homeomorfismo}
For every \(N\geq0\), the map
\begin{equation}
 \mathcal E_N:=\mathcal E^{\times N}\circ\beta_{{\rm blk},N}:
 Y_N\xrightarrow{\ \cong\ }
 (\mathscr L_{27}\times\mathscr L_{27})^N
 \label{corpus9:xii:rectas-nivel}
\end{equation}
is bijective. These bijections satisfy
\begin{equation}
 \operatorname{pref}_N\circ\mathcal E_{N+1}
 =\mathcal E_N\circ\rho^{\rm blk}_{N+1,N}
 \label{corpus9:xii:rectas-truncamiento}
\end{equation}
and determine the homeomorphism
\begin{equation}
 \mathcal E_\infty:=\mathcal E^{\mathbb N}\circ\beta_{\rm blk}:
 \Omega_\Gamma^{\rm cal}\xrightarrow{\ \cong\ }
 (\mathscr L_{27}\times\mathscr L_{27})^{\mathbb N}.
 \label{corpus9:xii:rectas-limite}
\end{equation}
Its inverse is
\(\beta_{\rm blk}^{-1}\circ(\mathcal E^{-1})^{\mathbb N}\).
\end{theorem}
\begin{proof}
Proposition~\ref{corpus9:xii:rectas-biyeccion} provides the effective
inverse of each factor \(\mathcal E\), using the line table and
coordinate interlacing. Consequently,
\(\mathcal E_N^{-1}=\beta_{{\rm blk},N}^{-1}
\circ(\mathcal E^{-1})^{\times N}\). For \(N=0\), both maps act on
the empty word. Coordinatewise application commutes with deletion of
the final factor. Using \eqref{eq:cobertura-729-ch} gives
\[
 \begin{aligned}
 \operatorname{pref}_N\mathcal E^{\times(N+1)}
           \beta_{{\rm blk},N+1}
 &=\mathcal E^{\times N}\operatorname{pref}_N
           \beta_{{\rm blk},N+1}\\
 &=\mathcal E^{\times N}\beta_{{\rm blk},N}
           \rho^{\rm blk}_{N+1,N}.
 \end{aligned}
\]
This is \eqref{corpus9:xii:rectas-truncamiento}. Applying the inverse
to each factor of a sequence of line pairs determines a unique word
sequence. The homeomorphism
\eqref{eq:homeomorfismo-calibrado-ch} reconstructs from it a unique
history in \(\Omega_\Gamma^{\rm cal}\). The compositions recover the
original history and the original sequence, respectively. In the
cylinder topologies, a basic open set fixes finitely many coordinates.
The map and its inverse transform these conditions into conditions on
the same number of coordinates; both are continuous. This proves the
topological assertion and compatibility of the limit with every finite
truncation.
\end{proof}

To retain the information in the complete state, let \(\zeta\) denote
its remaining records: frame, orientation, carry, route, memory, and
boundary data. On the already constructed admissible domain, the
enlarged transformation is
\[
 \widehat{\mathcal E}(w,\zeta)=(\mathcal E(w),\zeta).
\]
Its image domain retains the same compatibility conditions, expressed
through the explicit inverse of the first coordinate. If
\(F:\mathcal X\to\mathcal Y\) is an operator between state domains,
its transport is
\[
 F^{\mathcal E}
 =\widehat{\mathcal E}_{\mathcal Y}\circ F
       \circ\widehat{\mathcal E}_{\mathcal X}^{-1}.
\]
The identity \(F^{\mathcal E}\widehat{\mathcal E}_{\mathcal X}
=\widehat{\mathcal E}_{\mathcal Y}F\) follows by cancelling the two
inverse maps. For a relational prolongation
\(\mathcal R\subset\mathcal X\times\mathcal Y\), the corresponding
definition is
\[
 \mathcal R^{\mathcal E}
 =\{(\widehat{\mathcal E}_{\mathcal X}x,
       \widehat{\mathcal E}_{\mathcal Y}y):(x,y)\in\mathcal R\}.
\]
The inverse transformation recovers exactly \(\mathcal R\).
If the prolongation also retains a witness \(\omega\) in a domain
\(W\), its realization \(\widetilde{\mathcal R}\subset
\mathcal X\times W\times\mathcal Y\) is transported by
\[
 (x,\omega,y)\longmapsto
 (\widehat{\mathcal E}_{\mathcal X}x,\omega,
       \widehat{\mathcal E}_{\mathcal Y}y).
\]
The inverse acts on the two endpoints and retains \(\omega\). Thus
distinct witnesses remain distinct even when they share endpoints,
and the recorded multiplicity is preserved. Functionality is preserved
on the domains where it has already been established.

Coverage at arbitrary depth comes here from the domain
\(\Omega_\Gamma^{\rm cal}\) and its proved reader, whereas ordered
pairs of lines provide reversible coordinates for that domain.
The initial census of \(243\) words and the global capacity of
\(729\) prolongations per block retain their distinct roles.
The diagonal, incident, and skew classification accompanies each factor;
conditions selecting adjacent walks are transported as restrictions on
the complete sequence.
% END REV09_RECTAS27_PROLONGACION


The subcylinders are now defined as subsets of the limit, rather than as pairs
of suffixes. For \(x\in Y_N\) and \(b\in\mathbb F_3^6\), let
\begin{align}
 \mathscr Z_N(x)
   &:=\{\xi\in\Omega_\Gamma^{\rm cal}:
       \rho^{\rm blk}_{\infty,N}(\xi)=x\},
 \label{eq:cilindro-real-limite-ch}\\
 \mathscr Z_{N+1}(x;b)
   &:=\{\xi\in\Omega_\Gamma^{\rm cal}:
       \rho^{\rm blk}_{\infty,N+1}(\xi)
       =\Gamma_{b,1+54N}^{[54]}(x)\}.
 \label{eq:subcilindro-real-limite-ch}
\end{align}

\begin{theorem}[Effective partition into \(729\) refinement subcylinders]
\label{thm:particion-729-subcilindros-ch}
For every \(N\ge0\) and every \(x\in Y_N\), there is a disjoint union
\begin{equation}
 \mathscr Z_N(x)=
 \bigsqcup_{b\in\mathbb F_3^6}\mathscr Z_{N+1}(x;b).
 \label{eq:particion-real-729-ch}
\end{equation}
The \(3^6=729\) terms are nonempty clopen refinement subcylinders; each is
realized by fifty-four edges of the constructed TPK.
\end{theorem}

\begin{proof}
Let \(\xi\in\mathscr Z_N(x)\). Its level-\(N+1\) coordinate belongs to
\(Y_{N+1}\), truncates to \(x\), and, by the bijection
\(\beta_{{\rm blk},N+1}\), has a unique final block
\(b\in\mathbb F_3^6\). Identity
\eqref{eq:beta-concatenacion-demostrada-ch} then implies
\(\rho^{\rm blk}_{\infty,N+1}(\xi)
=\Gamma_{b,1+54N}^{[54]}(x)\); therefore, \(\xi\) belongs to one of the terms
on the right-hand side. The reverse inclusion follows immediately by
truncation. If two terms intersected, their level-\(N+1\) coordinates would
be equal and the ledger decoder would give \(b=b'\); hence the union is
disjoint.

Each endpoint \(\Gamma_b^{[54]}(x)\) admits an infinite tail obtained by
iterating the zero-block extension; consequently, every term is nonempty.
Since \(Y_{N+1}\) is finite and discrete, the terms are inverse images of
points under a continuous projection of the inverse limit and are therefore
clopen. The very definition of \(\Gamma_b^{[54]}\) establishes its
fifty-four edges.
\end{proof}

Finally, the collection \((\Phi_N)_N\) induces the homeomorphism of limits
\begin{equation}
 \Phi_\infty:\Omega_\Gamma^{\rm cal}
 \xrightarrow{\ \cong\ }\mathbb Z_3^6,\qquad
 \Phi_\infty(\xi):=
 \bigl(\Phi_N(\rho^{\rm blk}_{\infty,N}(\xi))\bigr)_{N\ge0},
 \label{eq:phi-infinito-hensel-ext-ch}
\end{equation}
where
\(\mathbb Z_3^6=\varprojlim_N(\mathbb Z/3^N\mathbb Z)^6\).
Equation \eqref{eq:entrelazamiento-hensel-ext-ch} guarantees that the
right-hand side is a compatible class, and the finite bijections prove
injectivity, surjectivity, and continuity in both directions. Under
\(\Phi_\infty\), \(\mathscr Z_N(x)\) is exactly the congruence class
\[
 \{z\in\mathbb Z_3^6:z\equiv\Phi_N(x)\pmod{3^N}\},
\]
and the sets \(\mathscr Z_{N+1}(x;b)\) are its \(729\) refinement classes
modulo \(3^{N+1}\). This is the effective junction between TPK extension,
limit cylinders, and the Hensel tower.

For \(b=(b_1,\ldots,b_6)\in\mathbb F_3^6\), set
\[
 \nu(b):=\sum_{i=1}^{6}[b_i]_3\,3^{6-i}
 \in\{0,\ldots,728\},
\tag{C.40c$_0$}
\label{eq:valor-bloque-ch}
\]
where \([b_i]_3\in\{0,1,2\}\) is the positional representative. The
\(243\) visible words in the initial census are the image of the first
emission; they are not to be confused with the \(3^6\) capacity of each
extension block. Finally,
\[
 X_{\infty,\mathbb Z}^{\rm cal}
 :=\mathbb Z_{\rm HMT}\times\Omega_\Gamma^{\rm cal}.
\]


For a prefix \(s=(m;b_0,\ldots,b_{N-1})\), define in the generated arithmetic
of \(\mathbb Q_{\rm HMT}\)
\[
 a_N(s):=m+\sum_{q<N}\frac{\nu(b_q)}{729^{q+1}},
 \qquad
 \operatorname{Succ}(s):=
 \{(m;b_0,\ldots,b_{N-1},b):b\in\mathbb F_3^6\}.
\]
If \((m;\beta_{{\rm blk},N}(x))\) denotes the prefix obtained by prepending
\(m\) to the tuple \(\beta_{{\rm blk},N}(x)\),
\eqref{eq:cobertura-729-ch} gives the typed equality
\[
 \operatorname{Succ}(m;\beta_{{\rm blk},N}(x))
 =\{(m;\beta_{{\rm blk},N+1}(y)):
       y\in\operatorname{Ext}^{\rm cal}_N(x)\}.
\]
From their first definition, the half-open cylinders are subsets of the HMT
completion:
\[
 I_N^{\rm HMT}(s):=
 \{r\in\mathbb R_{\rm HMT}:
      a_N(s)\le_{\rm HMT}r<_{\rm HMT}a_N(s)+729^{-N}\}.
\]
They satisfy materially
\[
 I_N^{\rm HMT}(s)=
 \bigsqcup_{y\in\operatorname{Succ}(s)}I_{N+1}^{\rm HMT}(y),
 \qquad |\operatorname{Succ}(s)|=729,
 \qquad \operatorname{diam}_{\rm HMT}I_N^{\rm HMT}(s)
      =729^{-N}\longrightarrow0.
\tag{C.40a}\label{eq:ch-particion-cilindros}
\]
The boundary convention identifies the two representations only after they
have been constructed. The partition~\eqref{eq:ch-particion-cilindros},
nesting, and contraction produce

\[
\begin{aligned}
 P_{\rm ar}&:X_{\infty,\mathbb Z}^{\rm cal}
     \twoheadrightarrow\mathbb R_{\rm HMT},\\
 P_{\rm ar}(m,\xi)
   &:=\text{the unique element of }
     \bigcap_N C_N^{\rm cl}
       \bigl(m;\beta_{{\rm blk},N}(\rho^{\rm blk}_{\infty,N}(\xi))\bigr),\\
 C_N^{\rm cl}(s)&:=\overline{I_N^{\rm HMT}(s)}.
\end{aligned}
\tag{C.40}\label{eq:ch-proyeccion-arquimediana}
\]

Therefore,

\[
 X_{\infty,\mathbb Z}^{\rm cal}/\!\sim_{P_{\rm ar}}
 \;\cong\;\mathbb R_{\rm HMT}.
\tag{C.41}\label{eq:ch-cociente-recta}
\]

The other representation, \(P_{\rm exc}\), preserves incidence, boundary, and
memory toward Paley--\allowbreak Witt--\allowbreak Golay--\allowbreak Leech and, from there, toward
\(V^\natural\)--Monster. Both arise from the same antecedent in
~\eqref{eq:ch-limite-inverso}; neither the Archimedean coordinate nor the
exceptional recognition selects the seeds. The five constructions
\(sp,sol,gau,coh,det\) remain consubstantial within
\(\mathcal C_{\rm cont}^{\rm disc}\) and are not disaggregated in order to
carry out the cardinal step.

\subsubsection{Descent of operations to the evaluation quotient}
\label{sec:ch-descenso-fibras}

The preceding quotient identifies histories publishing the same value, not
their complete registers. The following criterion characterizes exactly
when an operation on histories induces an operation on that quotient without
choosing a privileged history within each fiber; its formulation
continues the genealogical construction established in the preceding sections.

\begin{proposition}[Descent criterion through evaluation fibers]
\label{prop:ch-descenso-fibras}
Let \(q:\Omega\to J\) be the evaluation of a genealogical
chart onto its image. Let \(\star:D\to\Omega\) be an operation with
\(D\subseteq\Omega^2\) saturated by the fibers of \(q\times q\).
There exists a unique operation
\(\overline\star:(q\times q)(D)\to J\) such that
\[
 q(x\star y)=q(x)\,\overline\star\,q(y)
\]
if and only if, for pairs in the domain,
\[
 q(x)=q(x'),\quad q(y)=q(y')
 \quad\Longrightarrow\quad q(x\star y)=q(x'\star y').
\]
\end{proposition}

\begin{proof}
If the visible operation exists, its result depends only on the two
values, proving necessity. For sufficiency, define
\(s\,\overline\star\,t=q(x\star y)\), where \(q(x)=s\) and
\(q(y)=t\). Saturation preserves the domain when representatives are changed;
the stated condition preserves the result. Surjectivity of
\(q\times q\) onto the image of the domain proves uniqueness.
\end{proof}

A coordinate does not automatically replace a history as an argument of
all its operations: the criterion determines when that substitution is
legitimate. In particular, transporting an operation through a bijection
does not suffice to identify it with a native operation preceding the reading.
If a chart \(\Omega\) is compact and \(q\) is continuous and surjective
onto a Hausdorff space \(J\), evaluation also induces a
homeomorphism \(\Omega/{\sim_q}\to J\): it is a continuous bijection from a
compact space to a Hausdorff space. It is the \emph{quotient}, not necessarily the
space of histories, that is topologically identified with \(J\).


\subsubsection{Internal reconstruction of the chart and uniform lift}

Within \(\mathfrak G_{\rm HMT}(h,m_\infty)\), the reconstruction in
\eqref{eq:ch-limite-inverso}--\eqref{eq:ch-cociente-recta} is carried out with
the same codes and operators. Denote by
\(\mathbb R_{\rm HMT}^{\mathfrak G}\) the internal HMT completion. The
uniqueness of the complete Archimedean ordered field provides, \textbf{after}
its construction, the internal isomorphism
\[
 \eta:\mathbb R_{\rm HMT}^{\mathfrak G}
 \xrightarrow{\ \cong\ }\mathbb R^{\mathfrak G}.
\tag{C.41$_0$}\label{eq:ch-isomorfismo-interno}
\]
For the corresponding internal interval
\(I_N^{\mathfrak G}(s):=
 I_N^{\rm HMT}(s)\cap\mathbb R_{\rm HMT}^{\mathfrak G}\), the isomorphism
preserves endpoints and order:
\[
 \eta[I_N^{\mathfrak G}(s)]
 =\{r\in\mathbb R^{\mathfrak G}:
      a_N(s)\le r<a_N(s)+729^{-N}\}.
\tag{C.41$_0$a}\label{eq:ch-intervalos-internos}
\]
The surjectivity of~\eqref{eq:ch-proyeccion-arquimediana} covers every element
of that completion, not only those obtained in a finite execution.

So that “every” is not concealed in the definition of the codomain, we
construct the two injections relating the line to its internal power set. Let

\[
 E(A)=\sum_{n\ge0}\frac{2\mathbf1_A(n)}{3^{2n+2}}
 \qquad
 \bigl(A\in\mathcal P^{\mathfrak G}(\omega)\bigr).
\tag{C.41a}\label{eq:ch-incrustacion-cantor}
\]

This representation places \(0\) or \(2\) only in the even ternary positions;
its image is contained in \([0,1/4]\). It therefore never reaches the right
endpoint of the root cylinder, and two distinct subsets have a different first
even position: \(E\) is injective and is not subject to boundary ambiguity.
Write \(E_{\rm HMT}:=\eta^{-1}\circ E\).

Given \(A\), begin with the root cylinder containing \(E_{\rm HMT}(A)\). If
\(s_N(A)\) is the prefix already chosen, write
\(s_N(A)=(0;b_0(A),\ldots,b_{N-1}(A))\). The partition
~\eqref{eq:ch-particion-cilindros}, transported by \(\eta^{-1}\), supplies a
unique half-open subcylinder \(s_{N+1}(A)\) containing
\(E_{\rm HMT}(A)\). By~\eqref{eq:cobertura-729-ch}, that extension is also an
admissible TPK extension. By induction,
\[
 s_{N+1}(A)\!\upharpoonright N=s_N(A),\qquad
 E_{\rm HMT}(A)\in I_N^{\rm HMT}(s_N(A))\quad(N<\omega).
\tag{C.41b$_0$}\label{eq:ch-prefijos-uniformes}
\]
The recursion is a uniform formula with parameters \(A,h\). Since both are
internal, Replacement in the model of Lemma~\ref{thm:ch-6b} forms the complete
sequence of blocks \((b_q(A))_{q<\omega}\). Without confusing the tape with
the sequence of prefixes, define
\[
 \xi_A:=\Lambda\bigl((b_q(A))_{q<\omega}\bigr)
 \in\Omega_\Gamma^{\rm cal,\mathfrak G}.
\]
This yields the \textbf{uniform lift of the Cantor embedding}

\[
 Q:\mathcal P^{\mathfrak G}(\omega)
 \longrightarrow X_{\infty,\mathbb Z}^{\rm cal,\mathfrak G},
 \qquad
 Q(A):=(0,\xi_A),\qquad
 P_{\rm ar}(Q(A))=E_{\rm HMT}(A).
\tag{C.41b}\label{eq:ch-levantamiento-uniforme}
\]

The map \(Q\) is not called an inverse of \(P_{\rm ar}\): their composition
expresses the embedding \(E_{\rm HMT}\), not the identity on the entire line.
In the opposite direction, let \((q_n)_{n<\omega}\) be the generated
enumeration of \(\mathbb Q_{\rm HMT}\), and define
\[
 C(r):=\{n<\omega:\eta(q_n)<r\}
 \qquad(r\in\mathbb R^{\mathfrak G}).
\]
The rational cut belongs to \(\mathcal P^{\mathfrak G}(\omega)\), and \(C\)
is injective by density. The injections \(E\) and \(C\), followed by the
internal Cantor--Bernstein theorem, prove

\[
 P_{\rm ar}
 \bigl[X_{\infty,\mathbb Z}^{\rm cal,\mathfrak G}\bigr]
 =\mathbb R_{\rm HMT}^{\mathfrak G},
 \qquad
 \mathbb R_{\rm HMT}^{\mathfrak G}
 \cong\mathbb R^{\mathfrak G}
 \approx\mathcal P^{\mathfrak G}(\omega).
\tag{C.41c}\label{eq:ch-recta-potencia}
\]

Here an arbitrary real is used only as a universally quantified variable to
prove surjectivity of an already constructed tree. It selects neither APP,
TRIT, TPK, the constants, nor the transitions of that tree.

\subsection{Enumeration by genealogical codes and HMT condensation}

Rule~\eqref{eq:ch-operador-def} supplies derivation names from

\[
 (\ulcorner\varphi\urcorner,
  \sigma(p_0),\ldots,\sigma(p_{k-1})).
\tag{C.42}\label{eq:ch-codigo-nombres}
\]

For any object \(a\) of the hierarchy, let \(\beta(a)\) be its first stage and
\(\sigma(a)\) its least derivation tree in the recursive order of
Lemma~\ref{thm:ch-6b}. The term \(\sigma(a)\) is not required to be a natural
number when the stage is uncountable. The relation

\[
 a\prec_{\mathfrak G}b
 \quad\Longleftrightarrow\quad
 (\beta(a),\sigma(a))<_{\rm lex}
 (\beta(b),\sigma(b))
\tag{C.42a}\label{eq:ch-orden-genealogico}
\]

is the global genealogical well-order \(<_{\mathfrak G}\) induced by the
constructor. Repeated names are eliminated by taking the least one; neither a
real nor an enumeration of the line is inserted as data.

\begin{lemma}[Internal countability of levels below \(\omega_1^{\mathfrak G}\)]
\label{thm:ch-7}
For every \(\beta<\omega_1^{\mathfrak G}\), the level
\(\mathfrak G_\beta(h,m_\infty)\) is countable within
\(\mathfrak G_{\rm HMT}(h,m_\infty)\), and possesses a canonical surjection

\[
 e_\beta:\omega\twoheadrightarrow
 \mathfrak G_\beta(h,m_\infty).
\tag{C.43}\label{eq:ch-enumeracion-niveles}
\]

\end{lemma}
\begin{proof}
The initial level has a countable code. If \(e_\beta\) enumerates level
\(\beta\), the formula codes and finite tuples of parameters in
~\eqref{eq:ch-codigo-nombres} form a countable set and enumerate the successor.
If \(\lambda<\omega_1^{\mathfrak G}\) is a limit, there is internally a
surjection \(s:\omega\twoheadrightarrow\lambda\); take the
\(<_{\mathfrak G}\)-least one, in the order already constructed in
Lemma~\ref{thm:ch-6b}, and diagonalize the \(e_{s(n)}\). Finally, among all
surjections so obtained, define \(e_\beta\) as the
\(<_{\mathfrak G}\)-least. The recursion \(\beta\mapsto e_\beta\) is a
uniformly definable class of \(\mathfrak G\); every \(e_\beta\), and every
restriction of the recursion to an ordinal, belongs to \(\mathfrak G\). It
uses only \(h,u_{h,m}\) and the global well-order.
\end{proof}

\begin{lemma}[Condensation of internal reals below \(\omega_1^{\mathfrak G}\)]
\label{thm:ch-8}
Every real in \(\mathfrak G_{\rm HMT}(h,m_\infty)\) appears at a level whose
index is smaller than \(\omega_1^{\mathfrak G}\).

\end{lemma}
\begin{proof}
Let \(r\subseteq\omega\) be a real in \(\mathfrak G\), and let
\(B:=\mathfrak G_0(h,m_\infty)\). This base is transitive and internally
countable. With the codes of natural numbers and finite pairs already fixed,
\(\operatorname{Ord}\cap B=\omega+1\). By induction,
\[
 \operatorname{Ord}\cap\mathfrak G_\alpha=(\omega+1)+\alpha.
\]
The successor adds the ordinal formed by the ordinals of the preceding level;
it cannot add a larger one, because that would not be a subset of that level.
Limits are unions. Choose a limit \(\Theta\), with \(\Theta>\omega^2\), such
that \(r\in\mathfrak G_\Theta\). Thus
\(\operatorname{Ord}\cap\mathfrak G_\Theta=\Theta\), and every level
\(\mathfrak G_\xi\), with \(\xi<\Theta\), belongs to
\(\mathfrak G_\Theta\). Internally form the expanded structure
\[
 \mathcal A_\Theta=
 \langle\mathfrak G_\Theta,\in,u_{h,m},
         H_\Theta,<_{\mathfrak G}\!\upharpoonright\mathfrak G_\Theta\rangle
\]
with the level relation as a predicate of its signature,
\(H_\Theta=\{\langle\xi,x\rangle:\xi<\Theta\ \&\
x\in\mathfrak G_\xi\}\), and its definable Skolem functions. Let
\[
 Y=\operatorname{Hull}^{\mathcal A_\Theta}
   (B\cup\{B,r\})\prec\mathcal A_\Theta .
\]
Closure under the countable collection of Skolem functions shows within
\(\mathfrak G\) that \(Y\) is countable. Let \(c:Y\rightarrow M\) be its
Mostowski collapse, and let
\[
 \bar\Theta:=\operatorname{otp}(Y\cap\Theta)
 <\omega_1^{\mathfrak G}.
\]
Since \(B\subseteq Y\) is transitive, induction on membership shows that
\(c\) fixes every element of \(B\), as well as \(B\) itself. In particular,
it fixes \(\omega,u_{h,m},h,m_\infty\). Since
\(r\subseteq\omega\subseteq Y\), it also follows that \(c(r)=r\).

We prove by induction on \(\xi\in Y\cap\Theta\) that
\[
 c``\bigl(Y\cap\mathfrak G_\xi(h,m_\infty)\bigr)
 =\mathfrak G_{c(\xi)}(h,m_\infty).
\tag{C.43a}\label{eq:ch-condensacion}
\]
For \(\xi\in Y\cap\Theta\), the predicate \(H_\Theta\) defines the set
\(\mathfrak G_\xi\) from \(\xi\), so \(\mathfrak G_\xi\in Y\).
Relativizing each formula to that set gives
\(Y\cap\mathfrak G_\xi\prec\mathfrak G_\xi\), in the language into which
Lemma~\ref{thm:ch-6} translates the HMT definitions.

The base case of~\eqref{eq:ch-condensacion} has already been proved because the
collapse fixes \(B\). At a successor, if \(\xi+1\in Y\), then also
\(\xi\in Y\) and \(c(\xi+1)=c(\xi)+1\). For
\(a\in Y\cap\mathfrak G_{\xi+1}\), fix a finite formula defining it over
\(\mathfrak G_\xi\). The assertion
\[
 a=\{z\in\mathfrak G_\xi:
                  \varphi^{\mathfrak G_\xi}(z,\bar p)\}
\]
is an ordinary formula over \(\mathfrak G_\Theta\), with parameters
\(a,\mathfrak G_\xi\). Elementarity supplies the parameters \(\bar p\) within
\(Y\cap\mathfrak G_\xi\). The induction hypothesis and the elementarity of
that level show that \(c(a)\) is the subset defined by the same formula and
the collapsed parameters over \(\mathfrak G_{c(\xi)}\). Conversely, given a
formula and parameters at the collapsed level, lift the parameters to the
original level. The subset defined there exists, is unique, and belongs to
\(Y\) by elementarity. Its collapse is the required subset. This proves both
inclusions at the successor stage without introducing a uniform truth
predicate.

For a limit \(\lambda\in Y\cap\Theta\), elementarity supplies, for every
\(x\in Y\cap\mathfrak G_\lambda\), a witness \(\xi\in Y\cap\lambda\) with
\(H_\Theta(\xi,x)\). Therefore,
\[
 Y\cap\mathfrak G_\lambda
 =\bigcup_{\xi\in Y\cap\lambda}(Y\cap\mathfrak G_\xi),
 \qquad c``(Y\cap\lambda)=c(\lambda).
\]
Continuity of the hierarchy and the induction hypothesis give
~\eqref{eq:ch-condensacion} at the limit. Cofinality of \(Y\cap\lambda\) in
the ambient ordinal \(\lambda\) is not required: its collapsed indices range
over \(c(\lambda)\).

Finally, every \(x\in Y\) has a witness \(\xi\in Y\cap\Theta\) with
\(H_\Theta(\xi,x)\), and \(\bar\Theta\) is a limit. Taking unions,
\[
 M=\bigcup_{\eta<\bar\Theta}\mathfrak G_\eta(h,m_\infty)
  =\mathfrak G_{\bar\Theta}(h,m_\infty).
\]
Since \(c(r)=r\in M\), it follows that
\(r\in\mathfrak G_{\bar\Theta}\) with
\(\bar\Theta<\omega_1^{\mathfrak G}\).
\end{proof}

The lemma does not say that a finite-length prefix enumerates the line. It says
that every generated hereditary object has a first transfinite level and a
first genealogical code within that level. This is precisely the information
lost by the extensional reduct when it preserves only the expressed value.

\subsection{Injection of the generated line into the first uncountable ordinal of the genealogical closure}

Let \(\iota(r)\subseteq\omega\) be the rational cut of
\(r\in\mathbb R_{\rm HMT}^{\mathfrak G}\) relative to the enumeration of
\(\mathbb Q_{\rm HMT}\) previously constructed from APP integers and
quotients. The map \(r\mapsto\iota(r)\) is definable from \(r\) and that
enumeration, so \(\iota(r)\in\mathfrak G\); Lemma~\ref{thm:ch-8} places it at
a stage with countable index. Define

\[
 \begin{aligned}
  \beta(r)&:=\min\{\beta:
       \iota(r)\in\mathfrak G_{\beta+1}(h,m_\infty)\},\\
  n(r)&:=\min\{n:e_{\beta(r)+1}(n)=\iota(r)\},\\
  j_{\rm HMT}(r)&:=\omega\cdot\beta(r)+n(r).
 \end{aligned}
\tag{C.44}\label{eq:ch-indices-real}
\]

Definition~\eqref{eq:ch-indices-real} is uniform in the parameters
\(h,m_\infty\). Since \(\mathbb R_{\rm HMT}^{\mathfrak G}\) is a set of
\(\mathfrak G\), Separation fixes the domain of the definition, and
Replacement in \(\mathfrak G\) forms its graph:
\[
 \{\langle r,\omega\cdot\beta(r)+n(r)\rangle:
 r\in\mathbb R_{\rm HMT}^{\mathfrak G}\}\in\mathfrak G.
\]
Therefore, \(j_{\rm HMT}\) is an internal function of \(\mathfrak G\).

By Lemmas~\ref{thm:ch-7} and~\ref{thm:ch-8},
\(j_{\rm HMT}(r)<\omega_1^{\mathfrak G}\) for every generated real. If
\(j_{\rm HMT}(r)=j_{\rm HMT}(s)\), ordinal division by \(\omega\) recovers the
same pair \((\beta,n)\); \eqref{eq:ch-enumeracion-niveles} recovers the same
rational cut and, by completeness, \(r=s\). Hence

\[
 j_{\rm HMT}:\mathbb R_{\rm HMT}^{\mathfrak G}
 \hookrightarrow\omega_1^{\mathfrak G}
\tag{C.45}\label{eq:ch-inyeccion-j}
\]

is an injection constructed from the rank and name of the HMT genealogy. Its
existence uses precisely the semantic closure~\eqref{eq:ch-operador-def}, the
well-ordering of names from Lemma~\ref{thm:ch-6b}, and the condensation
~\eqref{eq:ch-condensacion}; it does not receive the Continuum Hypothesis in
order to choose the stage. The subsequent set-theoretic identification
describes this same machinery in different notation without reversing its
causal provenance.

\subsection{Injection of the first uncountable ordinal into the line and cardinal equivalence}

For every \(\beta<\omega_1^{\mathfrak G}\), Lemma~\ref{thm:ch-7} allows a subset
\(w_\beta\subseteq\omega\) to encode a countable well-order of type
\(\beta\). To avoid confusing the empty order with the
one-point order, both the domain \(D\subseteq\omega\) and its strict order
relation \(R\subseteq D^2\) are encoded. Once a bijective pairing
\(\langle\cdot,\cdot\rangle:\omega^2\to\omega\) is fixed, the code is
\[
 w(D,R)=\{2n:n\in D\}\cup
        \{2\langle a,b\rangle+1:(a,b)\in R\}.
\]
The even bits recover the domain, and the odd bits recover the relation. Take
the \(<_{\mathfrak G}\)-least real encoding such a well-order. Since the
property “encodes a well-order of type \(\beta\)” and the global well-order are
internal, Replacement forms the function \(\beta\mapsto w_\beta\) within
\(\mathfrak G\). Distinct order types produce distinct codes. The ternary map,
already used in~\eqref{eq:ch-incrustacion-cantor}, is defined by

\[
 E(w)=\sum_{n\ge0}\frac{2\mathbf1_w(n)}{3^{2n+2}}
\tag{C.46}\label{eq:ch-incrustacion-ternaria}
\]

and is injective and avoids the endpoint of the root cylinder. The lift
~\eqref{eq:ch-levantamiento-uniforme} ensures that
\(E_{\rm HMT}(w_\beta)\) belongs to the line expressed by an HMT history. Thus

\[
 k_{\rm HMT}(\beta):=E_{\rm HMT}(w_\beta),
 \qquad
 k_{\rm HMT}:\omega_1^{\mathfrak G}
 \hookrightarrow\mathbb R_{\rm HMT}^{\mathfrak G}.
\tag{C.47}\label{eq:ch-inyeccion-k}
\]

The injections~\eqref{eq:ch-inyeccion-j} and~\eqref{eq:ch-inyeccion-k},
constructed without using the equality sought, give by Cantor--Bernstein

\[
 \bigl|\mathbb R_{\rm HMT}^{\mathfrak G}\bigr|
 =\aleph_1^{\mathfrak G}.
\tag{C.48}\label{eq:ch-igualdad-cardinal}
\]

Since~\eqref{eq:ch-recta-potencia} identifies the complete HMT line with
\(\mathcal P^{\mathfrak G}(\omega)\), \eqref{eq:ch-igualdad-cardinal} rewrites
exactly as

\[
 \boxed{
  \mathfrak G_{\rm HMT}(h,m_\infty)
  \models 2^{\aleph_0}=\aleph_1.}
\tag{C.49}\label{eq:ch-decision}
\]

The equality does not arise from \(729\), from an informal product of
cardinality and ordinality, from the digits of a constant, or from a numerical
bracket. It arises from the covering of the line by the inverse limit, the
condensation of the genealogical closure, and the two explicit injections.

\subsection{Full quantifier over all internal subsets}

Let \(A\in\mathcal P^{\mathfrak G}
(\mathbb R_{\rm HMT}^{\mathfrak G})\). If \(A\) is countable, the first
alternative holds. If it is not, \(j_{\rm HMT}[A]\) is an uncountable subset
of \(\omega_1^{\mathfrak G}\). It cannot be bounded by a countable ordinal,
since every initial segment of \(\omega_1\) is countable. Its increasing
enumeration has order type \(\omega_1^{\mathfrak G}\), and the inverse of
\(j_{\rm HMT}\) on its image produces an injection
\(\omega_1^{\mathfrak G}\hookrightarrow A\). The inclusion
\(A\subseteq\mathbb R_{\rm HMT}^{\mathfrak G}\),
\eqref{eq:ch-igualdad-cardinal}, and Cantor--Bernstein give

\[
 \boxed{
 \forall A\in\mathcal P^{\mathfrak G}
   (\mathbb R_{\rm HMT}^{\mathfrak G}),\qquad
 |A|\le\aleph_0
 \quad\text{or}\quad
 |A|=\bigl|\mathbb R_{\rm HMT}^{\mathfrak G}\bigr|.}
\tag{C.50}\label{eq:ch-dicotomia}
\]

Here the full power set of the integral continuum generated by HMT is
quantified. The bijective transport of the folding with memory carries
~\eqref{eq:ch-dicotomia} to any of its representations without losing parts;
it does not replace the injections
\eqref{eq:ch-inyeccion-j}--\eqref{eq:ch-inyeccion-k}.

\subsection{Representation of the genealogical closure as \texorpdfstring{\(L[h,m_\infty]\)}{L[h,m-infinity]}}

Only after~\eqref{eq:ch-decision}--\eqref{eq:ch-dicotomia} is the construction
in~\eqref{eq:ch-clausura-transfinita} assigned its name in the usual
set-theoretic notation. Let \(L[u_{h,m}]\) be the constructible hierarchy
relative to the real already produced in~\eqref{eq:ch-codigo-conjunto}. Since
the HMT signature is countable and all its operators are graph relations
encoded by \(h\), Lemma~\ref{thm:ch-6} eliminates its symbols uniformly. This
yields the equality of classes

% The hyperlinked counter is separated from the editorial tag (51) to avoid
% two identically named PDF destinations in the chapter's last numbered equation.
\stepcounter{equation}
\[
 \mathfrak G_{\rm HMT}(h,m_\infty)
 =L[u_{h,m}]=L[h,m_\infty].
\tag{C.51}\label{eq:ch-representacion-L}
\]

To prove the first inclusion, strengthen the induction: for every ordinal
\(\beta\), both \(\mathfrak G_\beta\) and the sequence
\(\langle\mathfrak G_\xi:\xi\leq\beta\rangle\) belong to \(L[u_{h,m}]\).
The base is the transitive closure of a set constructed from \(\omega\) and
\(u_{h,m}\); the parameters \(h,m_\infty\) are recovered from \(u_{h,m}\).
At a successor, the preceding level is an \emph{internal set}, not merely a
subclass of \(L[u_{h,m}]\). Its satisfaction relation and the uniform
elimination~\eqref{eq:ch-eliminacion-uniforme} make it possible to form
internally the set of its definable subsets. At a limit, transfinite recursion,
Replacement, and Union in \(L[u_{h,m}]\) form the preceding sequence and its
union. Transitivity then gives \(\mathfrak G\subseteq L[u_{h,m}]\).

Conversely, Lemma~\ref{thm:ch-6b} supplies an internal transitive model of ZF,
containing all ordinals and \(u_{h,m}\) as an element. The recursion
constructing \(L[u_{h,m}]\) is carried out in that model: at every stage,
satisfaction over the preceding level and relative definability are absolute
between the two transitive domains. By induction, its levels coincide with
those of the ambient construction and belong to \(\mathfrak G\). Thus
\(L[u_{h,m}]\subseteq\mathfrak G\). Finally,
\eqref{eq:ch-codigo-conjunto} makes \(u_{h,m}\) and the pair
\((h,m_\infty)\) primitive-recursively interdefinable, from which the second
equality follows.

Equation~\eqref{eq:ch-representacion-L} explicitly identifies the closure with
the constructible hierarchy relative to the generated parameter. This
representation does not alter causal provenance: \(L[u]\) neither selects the
genealogical ledger nor defines the extensions or the line; the decision of
the Continuum Hypothesis follows from the exhaustivity rule and the maps proved
between \eqref{eq:ch-clausura-transfinita} and~\eqref{eq:ch-dicotomia}.
Taken together, the composition of \eqref{eq:ch-codigo-conjunto},
\eqref{eq:ch-clausura-transfinita}, \eqref{eq:ch-equivalencia-semantica}, and
\eqref{eq:ch-decision}--\eqref{eq:ch-representacion-L} and
Theorem~\ref{thm:decision-continuo-hmt} give a complete and uniform
expression of the realized universe, semantic coincidence, and cardinal
closure; none of these representations replaces APP--TRIT--TPK generation.

\subsection{Theorem giving an affirmative resolution of the Continuum Hypothesis}

\begin{theorem}[Decision of the continuum in HMT]
\label{thm:decision-continuo-hmt}
For every realized history
\(m_\infty\) produced by the joint action
\(\mathrm{APP}\to\mathrm{TRIT}\to\mathrm{TPK}\to
X_\infty^{\rm enr}\to\mathcal C_{\rm cont}^{\rm disc}\), the HMT genealogical
extension~\eqref{eq:ch-clausura-transfinita} generates the internal
reconstruction of the integral line
\eqref{eq:ch-proyeccion-arquimediana},
\(\mathbb R_{\rm HMT}^{\mathfrak G}\), constructs the injections
\eqref{eq:ch-inyeccion-j} and~\eqref{eq:ch-inyeccion-k}, and satisfies
\(2^{\aleph_0}=\aleph_1\) over its entire internal power set. Therefore, HMT
proves the Continuum Hypothesis affirmatively for the integral continuum that
it defines and generates. The notation \(L[h,m_\infty]\) is a subsequent
representation of that hierarchy, not the premise deciding the result.
\end{theorem}


\subsection{Falsification criteria for the cardinal closure}

The following list identifies failures that would invalidate a step of the
preceding closure. Construction, semantic, and cardinality controls are
distinguished:

\begin{enumerate}
\item there exists a level \(N\) or a surviving state without a compatible successor in
\eqref{eq:ch-extension-no-vacia};
\item \(u_{h,m}\) does not hereditarily encode \(h,m_\infty\), or some graph
relation of \(\Sigma_{\rm HMT}\) does not admit the uniform elimination~\eqref{eq:ch-eliminacion-uniforme};
\item an object admitted by the HMT semantic rules lacks a term in
\eqref{eq:ch-terminos}, or a term in~\eqref{eq:ch-terminos} expresses an object outside those rules, breaking
\eqref{eq:ch-equivalencia-semantica};
\item \(\mathfrak G\) does not satisfy one of the axioms used, its definable
well-order fails, or a purported enumeration \(e_\beta\) is not internal and
uniform;
\item the subcylinders in~\eqref{eq:ch-particion-cilindros} do not form an
exhaustive disjoint partition, their cylinders cease to contract, or the
recursion~\eqref{eq:ch-prefijos-uniformes} does not produce an internal
compatible history;
\item the collapse in Lemma~\ref{thm:ch-8} does not commute with a successor or
limit stage, or there exists an internal real that appears in no level with
index smaller than
\(\omega_1^{\mathfrak G}\);
\item the selection \(\beta\mapsto w_\beta\) is not internal or uniform, or
does not preserve order type;
\item two distinct reals have the same genealogical pair \((\beta,n)\) in
\eqref{eq:ch-indices-real}, two distinct ordinals produce the same code in
~\eqref{eq:ch-inyeccion-k}, or one of the injections
\(E,C,j_{\rm HMT},k_{\rm HMT}\) fails;
\item there exists an uncountable internal subset of the line whose range under
\(j_{\rm HMT}\) is bounded in \(\omega_1^{\mathfrak G}\);
\item some generative stage preceding the covering
~\eqref{eq:ch-particion-cilindros} uses the Continuum Hypothesis,
\(\mathbb R\), a target constant, the Committee on Data of the International
Science Council (CODATA), the Particle Data Group (PDG), or a subsequent
physical identification as a causal antecedent to define a seed, a relation
\(\mathscr L_\tau\), \(\operatorname{Ext}\), a coefficient, admissibility, or
an extension of the tree. The subsequent lift~\eqref{eq:ch-prefijos-uniformes}
does not constitute such an inversion: there an arbitrary already constructed
real is a universal variable and selects, within the previously proved
exhaustive partition, the unique subcylinder containing it.
\end{enumerate}

No execution restricted to a preset number of levels replaces the universal
quantifier of the inverse system; an effective failure of compatibility at
some depth would falsify it. No subsequent set-theoretic notation proves
~\eqref{eq:ch-decision} by itself. The closure jointly uses
Lemma~\ref{thm:ch-6b}, the line--power-set identification in
~\eqref{eq:ch-recta-potencia}, Lemmas~\ref{thm:ch-7}--\ref{thm:ch-8}, and the
two injections constructed from the genealogy.
